%% This is file `medima-template.tex',
%% 
%% Copyright 2018 Elsevier Ltd
%% 
%% This file is part of the 'Elsarticle Bundle'.
%% ---------------------------------------------
%% 
%% It may be distributed under the conditions of the LaTeX Project Public
%% License, either version 1.2 of this license or (at your option) any
%% later version.  The latest version of this license is in
%%    http://www.latex-project.org/lppl.txt
%% and version 1.2 or later is part of all distributions of LaTeX
%% version 1999/12/01 or later.
%% 
%% The list of all files belonging to the 'Elsarticle Bundle' is
%% given in the file `manifest.txt'.
%% 
%% Template article for Elsevier's document class `elsarticle'
%% with harvard style bibliographic references
%%
%% $Id: medima-template.tex 153 2018-12-01 11:38:32Z rishi $
%% $URL: http://lenova.river-valley.com/svn/elsarticle/trunk/medima-template.tex $
%%
%% Use the option review to obtain double line spacing
%\documentclass[times,review,preprint,authoryear]{elsarticle}

%% Use the options `twocolumn,final' to obtain the final layout
%% Use longtitle option to break abstract to multiple pages if overfull.
%% For Review pdf (With double line spacing)
%\documentclass[times,twocolumn,review]{elsarticle}
%% For abstracts longer than one page.
%\documentclass[times,twocolumn,review,longtitle]{elsarticle}
%% For Review pdf without preprint line
%\documentclass[times,twocolumn,review,nopreprintline]{elsarticle}
%% Final pdf
\documentclass[times,twocolumn,final]{elsarticle}
%%
%\documentclass[times,twocolumn,final,longtitle]{elsarticle}
%%


%% Stylefile to load MEDIMA template
\usepackage{medima}
\usepackage{framed,multirow}

%% The amssymb package provides various useful mathematical symbols
\usepackage{amssymb}
\usepackage{latexsym}

% Following three lines are needed for this document.
% If you are not loading colors or url, then these are
% not required.
\usepackage{url}
\usepackage{xcolor}
\usepackage{array}
\usepackage{arydshln}
\usepackage{subcaption}
\usepackage{hyperref}
\usepackage{makecell}
\usepackage{booktabs}
\usepackage{multirow}
\usepackage{caption}
\usepackage{bbding}
\usepackage{graphicx}
\usepackage{hyperref}
\usepackage{array}
\usepackage{arydshln}
\usepackage{subcaption}
\usepackage{hyperref}
\usepackage{makecell}
\usepackage{booktabs}
\usepackage{multirow}
\usepackage{amsmath}
\usepackage{caption}
\usepackage{graphicx}
\usepackage{bm}
\definecolor{newcolor}{rgb}{.8,.349,.1}

\journal{Medical Image Analysis}

\begin{document}

\verso{Given-name Surname \textit{et~al.}}

\begin{frontmatter}

\title{MedLSAM: Localize and Segment Anything Model for 3D CT Images}

\author[1,2]{Wenhui \snm{Lei}}
\author[3,4]{Wei \snm{Xu}}
\author[2,4]{Kang \snm{Li}}
\author[1,2]{Xiaofan \snm{Zhang}\corref{cor1}}
\author[1,2]{Shaoting \snm{Zhang}}
\cortext[cor1]{Corresponding author: xiaofan.zhang@sjtu.edu.cn}
\address[1]{School of Electronic Information and Electrical Engineering, Shanghai Jiao Tong University, Shanghai, China}
\address[2]{Shanghai AI Lab, Shanghai, China}
\address[3]{School of Biomedical Engineering, Division of Life Sciences and Medicine, University of Science and Technology of China, Hefei, China}
\address[4]{West China Biomedical Big Data Center, West China Hospital, Sichuan University, Chengdu, China}


\begin{abstract}
%%%
The Segment Anything Model (SAM) has recently emerged as a groundbreaking model in the field of image segmentation. Nevertheless, both the original SAM and its medical adaptations necessitate slice-by-slice annotations, which directly increase the annotation workload with the size of the dataset. We propose MedLSAM to address this issue, ensuring a constant annotation workload irrespective of dataset size and thereby simplifying the annotation process. Our model introduces a 3D localization foundation model capable of localizing any target anatomical part within the body. To achieve this, we develop a Localize Anything Model for 3D Medical Images (MedLAM), utilizing two self-supervision tasks: unified anatomical mapping (UAM) and multi-scale similarity (MSS) across a comprehensive dataset of 14,012 CT scans. We then establish a methodology for accurate segmentation by integrating MedLAM with SAM. By annotating several extreme points across three directions on a few templates, our model can autonomously identify the target anatomical region on all data scheduled for annotation. This allows our framework to generate a 2D bbox for every slice of the image, which is then leveraged by SAM to carry out segmentation. We carried out comprehensive experiments on two 3D datasets encompassing 38 distinct organs. Our findings are twofold: 1) MedLAM is capable of directly localizing any anatomical structure using just a few template scans, yet its performance surpasses that of fully supervised models; 2) MedLSAM not only aligns closely with the performance of SAM and its specialized medical adaptations with manual prompts but achieves this with minimal reliance on extreme point annotations across the entire dataset. Furthermore, MedLAM has the potential to be seamlessly integrated with future 3D SAM models, paving the way for enhanced performance. Our code is public at \href{https://github.com/openmedlab/MedLSAM}{https://github.com/openmedlab/MedLSAM}.
%%%%
\end{abstract}

\begin{keyword}
%% MSC codes here, in the form: \MSC code \sep code
%% or \MSC[2008] code \sep code (2000 is the default)
%% Keywords
\KWD SAM\sep Medical Image Segmentation\sep Contrastive Learning
\end{keyword}

\end{frontmatter}

%\linenumbers

%% main text
\section{Introduction}
\label{sec1}
\begin{figure*}[t!]
    \centering
    \includegraphics[width=1\linewidth]{seg_pipe.pdf}
    \caption{The overall segmentation pipeline of MedLSAM operates as follows. Given a dataset of any size, MedLSAM first applies a localization process (MedLAM) to identify the six extreme points (in the z, x, and y directions) of any anatomical structure of interest. This process results in the generation of a 3D bbox encompassing the targeted organ or structure. Subsequently, for each slice within this 3D bbox, a 2D bbox is derived, representing the projection of the 3D bbox onto that specific slice. These 2D bboxes are then utilized by the Segment Anything Model (SAM) to carry out precise segmentation of the target anatomy, thereby automating the entire segmentation process.}
    \label{fig:MedLSAM}
\end{figure*}

\par Recently, there has been an increasing interest in the field of computer vision to develop large-scale \textit{foundational models} that can concurrently address multiple visual tasks, such as image classification, object detection, and image segmentation. For instance, CLIP \citep{radford2021learning}, aligning a vast number of images and text pairs collected from the web, could recognize new visual categories using text \textit{prompts}. Similarly, GLIP \citep{li2022grounded, zhang2022glipv2}, which integrates object detection with phrase grounding in pre-training using web-sourced image-text pairs, boasts strong zero-shot and few-shot transferability in object-level detection tasks. On the segmentation front, the \textit{Segment Anything Model (SAM)} \citep{kirillov2023segment} has recently demonstrated remarkable capabilities in a broad range of segmentation tasks with appropriate prompts, e.g. points, bounding box (bbox) and text.

\par As foundational models increasingly demonstrate their prowess in general computer vision tasks, the medical imaging domain, characterized by limited image availability and high annotation costs, is taking keen notice of their potential. Such challenges in medical imaging underscore the pressing need for foundational models, drawing greater attention from researchers \citep{zhang2023challenges}. Some studies have endeavored to design self-supervised learning tasks tailored to dataset characteristics, pre-training on a vast amount of unlabeled data, and then fine-tuning on specific downstream tasks to achieve commendable results \citep{wang2023mis, zhou2023foundation, vorontsov2023virchow}. In contrast, other works have focused on pre-training with large annotated datasets, either to fine-tune on new tasks \citep{huang2023stu} or to enable direct segmentation without the necessity of further training or fine-tuning \citep{butoi2023universeg}. Overall, these models offer myriad advantages: they simplify the development process, reduce dependence on voluminously labeled datasets, and bolster patient data privacy.
\par 

% \par Advances in this direction have been made by models inspired by CLIP, such as ConVIRT \citep{zhang2022contrastive}, PubMedCLIP \citep{eslami2023pubmedclip}, and BiomedCLIP \citep{zhang2023large}. These models have primarily contributed to classification tasks by leveraging pre-trained vision-language grounding capabilities.

\par In the field of segmentation, models such as SAM and its medical adaptations \citep{MSA, MedSAM, cheng2023sam, he2023sam, huang2023segment, he2023accuracy, zhang2023customized, mazurowski2023segment, zhang2023segment} have demonstrated significant potential. Both \cite{huang2023segment} and \cite{cheng2023sam} have collected extensive medical image datasets and conducted comprehensive testing on SAM. They discovered that, while SAM exhibits remarkable performance on certain objects and modalities, it is less effective or even fails in other scenarios. This inconsistency arises from SAM's primary training on natural images. To enhance SAM's generalization in medical imaging, studies like \cite{MSA, MedSAM} have fine-tuned the original SAM using a large annotated medical image dataset, thereby validating its improved performance over the base SAM model in test datasets. Despite these advancements, the models still require manual prompts, such as labeled points or bounding boxes, which remains a substantial limitation due to the associated time and financial implications.
\par In scenarios requiring manual annotation prompts, annotating 3D images proves especially time-consuming and labor-intensive. A single 3D scan often contains dozens or even hundreds of slices. When each slice comprises multiple target structures, the annotation requirements for the entire 3D scan grow exponentially\citep{wang2018deepigeos, lei2019deepigeos}. To eliminate this vast annotation demand and facilitate the use of SAM and its medical variants on 3D data, an effective approach would be to automatically generate suitable prompts for the data awaiting annotation. While this could be achieved by training a detector specifically for the categories to be annotated, it introduces the additional burden of annotating data for detector training, and the detectable categories remain fixed\citep{baumgartner2021nndetection}. A more optimal solution would be to develop a universal localization model for 3D medical images. This model would only require users to provide minimal annotations to specify the category of interest, allowing for direct localization of the desired target in any dataset without further training. 
\par However, to the best of our knowledge, foundational models specifically tailored for 3D medical image localization are still limited. In the domain of 2D medical image localization, MIU-VL \citep{qin2022medical} stands out as a pioneering effort. This approach ingeniously marries pre-trained vision-language models with medical text prompts to facilitate object detection in medical imaging. However, the application of these pre-trained models to specialized medical imaging domains, particularly 3D modalities like CT scans, remains challenging. This limitation is attributable to a significant domain gap, as these models were originally trained on 2D natural images.
\par \textbf{In response to the lack of localization foundation models for 3D medical images and the persistent need for extensive interaction in SAM, we introduce MedLSAM in this paper}. MedLSAM is an automated medical image segmentation solution that combines the Localize Anything Model for 3D Medical Image (MedLAM), a dedicated 3D localization foundation module designed for localizing any anatomical structure using a few prompts, with existing segmentation foundation models. As illustrated in Fig. \ref{fig:MedLSAM}, MedLSAM employs a two-stage methodology. The first stage involves MedLAM to automatically identify the positions of target structures within volumetric medical images. In the subsequent stage, the SAM model uses the bboxes from the first stage to achieve precise image segmentation. The result is a fully autonomous pipeline that eliminates the need for manual intervention.

The localization foundation model, MedLAM, is an extension of our conference version \citep{lei2021contrastive} and is premised on the observation that the spatial distribution of organs maintains strong similarities across different individuals. And we assume that there exists a standard anatomical coordinate system, in which the same anatomy part in different people shares similar coordinates. Therefore, we can localize the targeted anatomy structure with a similar coordinate in the unannotated scan. In our previous version, the model was trained separately for different anatomical structures, each using only a few dozen scan images. In contrast, this current study significantly expands the dataset to include 14,012 CT scans from 16 different datasets. This allows us to train a unified, comprehensive model capable of localizing any structure across the entire body. The training process involves a projection network that predicts the 3D physical offsets between any two patches within the same image, thereby mapping every part of the scans onto a shared 3D latent coordinate system.

% While our localization approach is robust, we acknowledge that individual variations in anatomical positioning could result in different anatomical structures sharing the same latent coordinates across various images. To mitigate this issue, we refine our localization accuracy by extracting pixel-level features from our points of interest. This method enables us to identify the most similar feature within the vicinity of the initially localized point, thus enhancing the overall localization accuracy. Our approach is inspired by the self-supervised learning tasks proposed in \cite{yao2021one}, which strive to maximize the similarity between the original and augmented instances of the same image point and minimize the similarity between different points. This technique can further improve the accuracy and precision of our model's localization.

For the segmentation stage, we employ the original SAM and the well-established MedSAM \citep{MedSAM} as the foundation for our segmentation process. MedSAM, previously fine-tuned on a comprehensive collection of medical image datasets, has exhibited considerable performance in 2D and 3D medical image segmentation tasks. The use of such a robust model bolsters the reliability and effectiveness of our proposed pipeline.

Our contributions can be summarized as follows:
\par - We introduce MedLSAM, the first completely automated medical adaptation of the SAM model. It not only eliminates the need for manual intervention but also propels the development of fully automated foundational segmentation models in medical image analysis. To the best of our knowledge, this is the first work to achieve complete automation in medical image segmentation by integrating the SAM model;

\par - Addressing the gap in localization foundation models for 3D medical imaging, we present MedLAM. This model, an extension of our prior conference work \citep{lei2021contrastive}, incorporates pixel-level self-supervised tasks and is trained on a comprehensive dataset of 14,012 CT scans covering the entire body. As a result, it's capable of directly localizing any anatomical structure within the human body and surpassing the fully-supervised model by a significant margin;

\par - We validate the effectiveness of MedLSAM on two 3D datasets that span 38 organs. As showcased in Table \ref{tab:2}, our localization foundation model, MedLAM, exhibits localization accuracy that not only significantly outperforms existing medical localization foundational models but also rivals the performance of fully-supervised models. Concurrently, as demonstrated in Table \ref{tab:4}, MedLSAM aligns with the performance of SAM and its medical adaptations, substantially reducing the reliance on manual annotations. To further the research in foundational models for medical imaging, we have made all our models and codes publicly available.

\section{Related Work}
\subsection{Medical Visual Foundation Models}
%过去visual大模型进展，SAM出现，在医学领域的使用。评测下来bbox效果最好
% Recent breakthroughs in large language models (LLM), including ChatGPT \citep{ouyang2022training} and LLaMA \citep{touvron2023llama}, have consistently demonstrated impressive zero-shot and few-shot capabilities across various NLP tasks. Their potent performance largely stems from the vast text datasets sourced from the web and the scalable model architecture, primarily the Transformer \citep{vaswani2017attention}. Analogously, foundational models in the visual domain have benefited from the combination of expansive data and the increasing scaling of the Transformer structure.  
% \par One prominent example is CLIP (Contrastive Language-Image Pretraining) \citep{radford2021learning}, which has revolutionized the multi-modality learning. By training on a large dataset of image-text pairs, CLIP enables zero-shot transfer learning across a wide range of visual tasks. It aligns the feature space of images and texts so effectively that a model trained in this manner can understand and process new visual categories based merely on textual descriptions. In light of this context, ConVIRT \citep{zhang2022contrastive} has pioneered the enhancement of image representation capabilities by maximizing the agreement between true image-text pairs, subsequently achieving promising transfer learning results in four medical image classification tasks. Drawing inspiration from CLIP, PubMedCLIP \citep{eslami2023pubmedclip} and BiomedCLIP \citep{zhang2023large} have employed the Medical Visual Question Answering (VQA) dataset and figure-caption pairs extracted from biomedical research articles in PubMed Central, respectively, to fine-tune specialized CLIP models tailored for the medical image domain. 
% \par For object detection, MIU-VL\cite{qin2022medical} first utilized pre-trained vision-language models such as GLIP\cite{li2022grounded}, combined with auto-generated medical text prompts, to effectively transfer domain knowledge from natural to medical images. While this approach yielded promising results, the inherent limitations arise from the model being trained on natural images and its 2D nature. In our experiments, we observed that it struggled to proficiently handle 3D scanning data, such as CT scans. 
\par Over recent years, large-scale pre-trained deep learning vision models have marked remarkable milestones with models like ResNet\citep{he2016deep}, ConvNeXt\citep{liu2022convnet}, Vision Transformers (ViT)\citep{dosovitskiy2020image}, and DINO\citep{caron2021emerging}. Owing to their profound representational power and exceptional generalization capabilities, these models have consistently excelled in various downstream applications.
\par In the realm of medical image analysis, the challenges are compounded by concerns of data privacy and the often tedious nature of annotation tasks, making the procurement of large-scale annotated data both challenging and costly. This elevates the significance of visual foundational models in the field. For instance, REMEDIS\citep{azizi2022robust} harnesses the potential of self-supervised learning on a diverse set of medical images spanning multiple modalities. Impressively, it rivals the performance of heavily supervised baselines in downstream disease classification tasks, requiring only between 1\% to 33\% of the typical retraining data. Meanwhile, SAIS\citep{azizi2022robust}, capitalizing on the foundational model DINO\citep{caron2021emerging} pre-trained on standard images, showcases impressive generalization and prowess in specialized tasks, such as surgical gesture recognition and skill evaluation.
\par Turning the attention to segmentation, a key aspect within medical imaging, there has been a surge of interest in the application of SAM for medical image segmentation \citep{cheng2023sam, he2023sam, MSA, MedSAM, huang2023segment, he2023accuracy, zhang2023customized, mazurowski2023segment, zhang2023segment}. Notably, with additional medical data collected for fine-tuning, the medical adaptations tailored for SAM, specifically \cite{MSA, MedSAM}, have exhibited commendable results in this domain.

However, a significant hurdle for SAM's utility in medical image segmentation is the necessity for manual annotations. This involves annotating specific points or bboxes that demarcate the segmentation region, a process that is both time-consuming and expensive. Intriguingly, comprehensive studies by \cite{huang2023segment} and \cite{cheng2023sam} have delved into the optimal prompts for SAM's application in medical imaging. They primarily investigated prompt types including bboxes, a combination of bboxes with foreground and background points, and the exclusive use of foreground and background points. Notably, in the majority of tasks, bboxes alone showcased superior performance.

In light of these findings, one of the central focuses of this paper is the automation of generating bbox prompts for SAM. Furthermore, we aim to advance the full automation of SAM's application in the medical domain, reducing the reliance on manual annotations and propelling its utility.

\begin{figure*}[t!]
\centering
\includegraphics[width=0.85\linewidth]{train.pdf}
\caption{The learning process of MedLAM. }
\label{fig:training}
\end{figure*}

\subsection{Medical Localization Models}
Traditionally, localization tasks, such as object detection, have relied on extensive annotations for a limited set of categories to enable models to produce corresponding bboxes \citep{redmon2016you, he2017mask, dai2016r, ren2015faster}. This challenge has also been extensively addressed in the realm of medical imaging, with a focus on areas such as organ localization \citep{xu2019efficient,navarro2020deep, hussain2021cascaded, navarro2022unified}, keypoint detection\citep{chen2021fast, wan2023multi}, and lesion detection\citep{yan2018deeplesion, zhang2018automatic}. Both organ localization and keypoint detection often serve as preliminary steps for many medical image analysis tasks, including image registration\citep{sun2013validation}, organ segmentation\citep{liu2023clip}, and disease diagnosis\citep{bilic2023liver}. While previous fully-supervised models have achieved commendable results in predefined tasks, their primary limitation is their difficulty in generalizing to new categories. The high cost and challenges associated with data collection and annotation further exacerbate this issue, especially pronounced in the field of medical imaging.
\par To reduce annotation costs, our previous work, the Relative Position Regression-Net (RPR-Net) \citep{lei2021contrastive}, along with another study, Cascade Comparing to Detect (CC2D) \citep{yao2021one}, pioneered the localization of regions of interest or landmarks using annotations from just a single instance. Specifically, CC2D learns to project pixels associated with the same anatomical structures into embeddings with high cosine similarities by addressing a self-supervised "patch matching" proxy task in a coarse-to-fine manner. However, both RPR-Net and CC2D were trained on limited datasets as independent models, without evolving into foundational models for localization. 

\section{Method}

As illustrated in Fig. \ref{fig:MedLSAM}, MedLSAM comprises two core components: the automated localization algorithm MedLAM and the automated segmentation algorithm SAM. In the subsequent sections, we first delve into the architecture and training of MedLAM. Following that, we detail how MedLAM synergizes with SAM to realize the complete autonomous segmentation framework.

\subsection{Training of MedLAM}
\label{sec:2.2}

Our MedLAM model, depicted in Fig. \ref{fig:training}, consists of three primary modules: a feature encoder, a multilayer perceptron (MLP), and a feature decoder. Both the feature encoder and decoder are composed of four convolutional blocks. The encoder blocks each have two convolutional layers followed by a downsampling layer, while the decoder blocks contain two convolutional layers succeeded by an upsampling layer. The MLP is structured with three fully connected layers.

\par We start by selecting a volumetric image $\bm{v}$ from the unannotated training set. We then extract two large image patches from $\bm{v}$, which serve as the source for further extractions. These large patches undergo a variety of transformations to produce two pairs of patches, namely, the original patch pair $\bm{x}_q, \bm{x}_s \in R^{D \times W \times H}$ and the augmented patch pair $\bm{x_q'}, \bm{x_s'}\in R^{D \times W \times H}$. The augmentation techniques employed include rotation, elastic transformations, and the addition of random noise. Our training objectives are twofold:
1) \textbf{Unified Anatomical Mapping (UAM)}: To project scan images from different individuals onto a shared anatomical coordinate space by predicting the relative distance between the original image patches $\bm{x}_q$ and $\bm{x}_s$.
2) \textbf{Multi-Scale Similarity (MSS)}: Ensure that the similarity between features corresponding to the same region is maximized in both the original image patches $\bm{x}_q, \bm{x}_s$ and their augmented counterparts $\bm{x_q'}, \bm{x_s'}$.


\subsubsection{Unified Anatomical Mapping (UAM)}
\label{sec:3.1.1}
In this step, we introduce the Unified Anatomical Mapping (UAM) framework, building upon the foundation of our earlier methodology \citep{lei2021contrastive}. The central principle remains consistent: mapping 3D scan images from diverse individuals onto a unified implicit 3D anatomical coordinate system. By doing so, identical anatomical structures from various individuals can be associated with the same coordinate, facilitating an initial, coarse localization of the point within the unannotated scan corresponding to our point of interest.

In contrast to our prior endeavors, UAM embodies a broader vision of creating a model that captures the underlying anatomical structures across a wide range of human body scans. While our previous methods primarily focused on training specialized models for individual localized scan datasets, the UAM framework takes advantage of an extensive collection of CT scans, covering the entire human spectrum. This approach aims to provide a more holistic model capable of operating seamlessly across different organs and body regions, increasing the versatility and applicability of our method but also ensuring higher accuracy and reliability in diverse medical contexts.

The RDR model aims to predict the 3D offset between the query patch $\bm{x}_q$ and the support patch $\bm{x}_s$. Considering $\bm{e} \in R^3$ as the pixel spacing of $\bm{v}$, and $\bm{c}_q, \bm{c}_s \in R^3$ as the centroid coordinates of $\bm{x}_q$ and $\bm{x}_s$ in $\bm{v}$ respectively, the ground truth offset $\bm{d}'_{qs}$ from $\bm{x}_q$ to $\bm{x}_s$ in the physical space can be calculated as:

\begin{equation}
\bm{d}'_{qs}= (\bm{c}_s-\bm{c}_q) \cdot \bm{e} .
\end{equation}

Both $\bm{x}_s$ and $\bm{x}_q$ undergo processing via the feature encoder to distill high-level features. Subsequently, MLP maps these features to their corresponding 3D latent vectors, $\bm{p}_s$ and $\bm{p}_q \in R^3$. This process leads to the predicted offset $\bm{d}_{qs} \in R^3$ from the query patch $\bm{x}_q$ to the support patch $\bm{x}_s$ being computed as:

\begin{equation}
 \bm{d}_{qs}= r \cdot tanh(\bm{p}_{s}-\bm{p}_{q}) .
\end{equation}

The utilization of the hyperbolic tangent function $tanh$ in conjunction with the hyper-parameter $r$ is intended to dictate the upper and lower bound of $\bm{d}_{qs}$, thereby covering the largest feasible offset. Lastly, to measure the difference between $\bm{d}_{qs}$ and $\bm{d}'_{qs}$, we employ the Mean Square Error (MSE) loss function:

\begin{equation}
L_{MSE} = ||\bm{d}_{qs}- \bm{d}'_{qs}||^2 .
\end{equation}

Owing to the constraints set by the hyper-parameter $r$ and the inherent anatomical similarities across different individuals, the model effectively projects the same anatomical regions from various individuals onto similar implicit coordinates after training. This projection mechanism aids in achieving a consistent and unified representation across diverse datasets, thereby enhancing the adaptability and precision of our MedLAM model.

\begin{figure*}[ht!]
    \centering
    \includegraphics[width=0.8\linewidth]{MSS.pdf}
    \caption{Deatils of the Multi Scale Similarity (MSS).}
    \label{fig:mss}
\end{figure*}

\subsubsection{Multi Scale Similarity (MSS)} 
However, given the inherent variations in anatomical positioning across different individuals, regions sharing the same latent coordinates in various images may still correspond to different anatomical structures. Therefore, we need to further refine the precision of our localization by extracting local pixel-level features from our points of interest. This allows us to pinpoint the most similar feature within the vicinity of the initially localized point, thereby enhancing the overall localization accuracy. This is inspired by the work in Cascade Comparing to Detect (CC2D) \citep{yao2021one}, which ensures that augmented instances of the same image yield highly similar features for the same point, while different points exhibit substantially divergent features.
\par More specifically, as illustrated in Fig. \ref{fig:mss}, the inputs to our MSS process encompass normalized multi-scale feature maps, denoted as $\bm{F}=\{\bm{F}^i,i=0,1,2\}$ and $\bm{F}'=\{\bm{F}^{'i},i=0,1,2\}$, extracted from $\bm{x}$ and $\bm{x}'$, respectively. These multi-scale feature maps consist of three different scales, where the size of each scale is $1/{2^i}$ times the size of the original image, with the largest scale being the same size as the original image.

\par Given a chosen point $\bm{c}_j$ from $\bm{x}$, which has a corresponding point in $\bm{x}'$ is denoted as $\bm{c}'_j$. For every scale $i$ of the multi-scale feature maps, we extract the feature vectors $\bm{f}_{\bm{c}_j}^i$ and $\bm{f}_{\bm{c}'_j}^{'i}$ corresponding to points $\bm{c}_j$ and $\bm{c}'_j$ respectively. Formally, $\bm{f}_{\bm{c}_j}^i=\bm{F}^i(\bm{c}_j)$ and $\bm{f}_{\bm{c}'_j}^{'i}=\bm{F}^{'i}(\bm{c}'_j)$.

\par To compute the similarity between these feature vectors and the corresponding scale feature maps in $\bm{x'}$, we use the cosine similarity metric. The cosine similarity $s$ between two vectors $\bm{a}$ and $\bm{b}$ is given by:

\begin{equation}
s(\bm{a}, \bm{b}) = \frac{\bm{a} \cdot \bm{b}}{||\bm{a}||_2 \cdot ||\bm{b}||_2} .
\end{equation}

\par Since the $\bm{F}^i$ and $\bm{F}'^i$ are both normalized, for each scale $i$, the cosine similarity map is computed by using the feature vector $\bm{f}_{\bm{c}_j}^i$ as a $1\times1\times1$ convolutional kernel that operates on the corresponding feature map from the augmented image. Our optimization objective is to increase the similarity between $\bm{f}_{\bm{c}_j}^i$ and $\bm{f}_{\bm{c}'_j}^{'i}$ relative to other points, effectively enhancing their mutual likeness while diminishing their similarity to the remaining features in the augmented image. To achieve this, we apply the softmax operation across all values in the cosine similarity map, resulting in the final probability map $\bm{S}^i_{\bm{c}_j}$. Correspondingly, the ground truth $\bm{Y}^i_{\bm{c}_j}$ at scale $i$, is set as: 

\begin{equation}
\bm{Y}^i_{\bm{c}_j}(\bm{c}) = 
\begin{cases} 
1, & \text{if}\ \bm{c} = round(\frac{\bm{c}'_j}{2^i}) \\
0, & \text{otherwise}
\end{cases}
\end{equation}
This ensures that all other points are set to 0, while the corresponding position for $\bm{c}'_j$ at scale $i$ is set to 1. We utilize the Cross Entropy (CE) loss to align each $\bm{S}^i_{\bm{c}_j}$ with its corresponding ground truth $\bm{Y}^i_{\bm{c}_j}$:

\begin{equation}
    L_{CE} = \sum_{i} -\bm{Y}^i_{\bm{c}_j}\cdot log(\bm{S}^i_{\bm{c}_j}) - (1-\bm{Y}^i)\cdot log(1-\bm{S}^i_{\bm{c}_j}).
\end{equation}
Eventually, we set the final loss function as:

\begin{equation}
    L_{SSL} = L_{MSE} +L_{CE}.
\end{equation}


%After resizing the resulting similarity maps to the original image size, we aggregate them. This process allows us to pinpoint the location within $\bm{x'}_s$ that exhibits the highest similarity to point $c_1$, thereby further refining our localization.

\begin{figure}[t!]
    \centering
    \includegraphics[width=1\linewidth]{localize_pipe.pdf}
    \caption{Structure of the inference stage in the Medical Localization Anything Model (MedLAM). The process involves moving an agent from a randomly initialized position toward the target landmark, guided by 3D relative distance and multi-scale feature vectors.}
    \label{fig:inference}
\end{figure} 

\begin{figure}[t!]
    \centering
    \includegraphics[width=1\linewidth]{sub_patch.pdf}
    \caption{Comparison between Whole-Patch Localization (WPL) and Sub-Patch Localization (SPL) strategies. (a) WPL generates a single minimal 3D bounding box ($\bm{B}$) by identifying the six extreme points of the target anatomical structure. (b) SPL divides the target structure into multiple segments, each with its own localized 3D bounding box ($\bm{B}_i$), resulting in a more accurate representation at each slice.}
    \label{fig:subpatch}
\end{figure} 

\subsection{Inference of MedLSAM}
\label{sec:3.2}
\par  In this section, we will elaborate on the process of utilizing MedLAM to localize any point of interest within scans. Additionally, we further elucidate the mechanism through which MedLAM crafts 3D bbox prompts for the anatomical structures of interest. These 3D prompts are then sliced into 2D bboxes corresponding to each layer of the structure. Each 2D bbox, acting as a prompt for its respective layer, is combined with the corresponding image slice and processed by the SAM model to produce the segmentation results.

\subsubsection{Landmark Localization}
\par As posited in Section \ref{sec:3.1.1}, we assume that the same anatomical structure in different individuals corresponds to the same latent coordinates. This allows us to calculate the latent coordinates of any point within the query volume and subsequently determine its relative distance to the target landmark by computing the difference in their latent coordinates. 
\par As shown in Fig. \ref{fig:inference}, we conceptualize the localization task as moving an agent from a randomly initialized position in the query image toward the target location. A patch is extracted from a random position within the query image, and simultaneously, a support patch is extracted from the support image, centered around the pre-specified landmark. Upon processing these two patches through the MedLAM model, we obtain a 3D offset that represents the estimated relative spatial displacement between the start and target positions and the multi-scale feature vectors of the pre-specified landmark. By updating the agent's location based on this offset, we achieve a coarse localization of the landmark within the query image. 

\par To further refine the landmark localization, we employ the Multi-Scale Similarity (MSS) component of MedLAM. Specifically, we extract a 3D patch of size $D\times W\times H$ centered around the coarsely localized point in the query image. Within the extracted patch, we compute multi-scale feature maps, perform similarity calculations, and aggregate the resulting similarity maps. This allows us to identify the location within the patch that exhibits the highest feature similarity, thereby significantly enhancing the precision of our landmark localization.

\subsubsection{Sub-Patch Localization (SPL)}
\par  As shown in Fig. \ref{fig:subpatch}(a), for a target anatomical structure, we can generate a single minimal 3D bounding box (bbox) $\bm{B}$ that completely encompasses the organ by locating its six extreme points in three dimensions:  $[\bm{z}_{min}, \bm{z}_{max}, \bm{x}_{min}, \bm{x}_{max}, \bm{y}_{min}, \bm{y}_{max}]$. We call this localization approach as \textbf{Whole-Patch Localization (WPL)}. However, the cross-sectional area of this 3D bbox on each layer may significantly exceed the actual cross-sectional area of the organ at that layer, resulting in highly inaccurate bbox prompts.
\par In response, we introduce the \textbf{Sub-Patch Localization (SPL)} strategy. Rather than using a single bounding box, SPL divides the target structure into multiple segments, each segment having its own dedicated 3D bbox, denoted as $\bm{B}_i$. This segmented approach, as visualized in Fig. \ref{fig:subpatch}(b), ensures that each bbox is tailored to more accurately represent the organ's actual cross-sectional area at every layer.
\par Delving into the specifics, for any given target organ, SPL slices the organ into segments at intervals of $n$ mm, resulting in $m$ distinct segments for the organ. Consequently, MedLAM is then tasked with localizing these $m$ segments in the query volume, producing $m$ individual bounding boxes: $[\bm{z}^i_{min}, \bm{z}^i_{max}, \bm{x}^i_{min}, \bm{x}^i_{max}, \bm{y}^i_{min}, \bm{y}^i_{max}]$ for each segment $i$. This refined approach provides a much more accurate spatial representation of the anatomical structure across its entirety.

\subsubsection{SAM Segmentation}
\par After generating the bbox prompts for each slice, we transition to the segmentation stage. For this, we utilize both SAM and MedSAM, a specialized variant of SAM that has been fine-tuned for medical image datasets. Both models serve as the foundation for our segmentation tasks. The versatility of SAM and the domain-specific adaptations of MedSAM help us provide robust segmentation results, thereby adding to the overall efficacy of the MedSLAM system. 



\section{Experiments}

\subsection{Dataset}
\par Our MedLAM model is trained on an extensive set of 16 datasets, which collectively comprise a total of 14,012 CT scans. These scans encompass various regions of the human body, providing comprehensive anatomical coverage. An overview of the training datasets is provided in Table~\ref{tab:dataset}. The diverse and abundant training data ensures the robustness and generalizability of our model across different medical imaging contexts.
\par To validate the effectiveness of our approach, we integrate MedLAM with two segmentation backbones: SAM\cite{kirillov2023segment} and MedSAM\cite{MedSAM}. We test these combined models on two CT segmentation datasets: 1) StructSeg19 Task1 dataset \footnote{https://structseg2019.grand-challenge.org/Home/} for the 22 Head-and-Neck (HaN) organs with 50 scans; 2) the WORD dataset \cite{luo2022word} for the 16 abdomen organs with 120 scans. 

\begin{table}[t!]
\caption{Detailed information of the 16 CT datasets for MedLAM training.}
\centering
\scalebox{0.75}{
\begin{tabular}{l|cc} 
\hline
Dataset          & Number       & Anatomical Region                    \\ \hline
GLIA \cite{bo2021toward}            & 1338         & HaN                           \\
ACRIN 6685\cite{lowe2019multicenter}       & 260          & HaN                           \\
OPC-Radiomics\cite{kwan2018radiomic}    & 606          & HaN                           \\
Head-Neck-PET-CT\cite{vallieres2017radiomics} & 298          & HaN                          \\
HNSCC \cite{grossberg2018imaging}           & 591          & HaN/Thorax/Abdomen     \\
autoPET \cite{gatidis2022whole}         & 1014         & Whole                       \\
MELA \cite{mela}          & 770          & Thorax                      \\
LIDC-IDRI\cite{armato2011lung}        & 1308         & Thorax                       \\
STOIC2021\cite{revel2021study}        & 2000         & Thorax                       \\
MSD-Lung\cite{antonelli2022medical}         & 95           & Thorax                      \\
CBIS-DDSM\cite{lee2017curated}        & 2620         & Thorax                      \\
AMOS 2022\cite{ji2022amos}        & 500          & Thorax/Abdomen             \\
Kits19\cite{heller2020state}           & 141          & Abdomen                     \\
MSD-Colon\cite{antonelli2022medical}        & 190          & Abdomen                     \\
MSD-Pancreas\cite{antonelli2022medical}     & 281          & Abdomen                      \\
FLARE2022\cite{ma2022fast}        & 2000         & Abdomen                      \\ \hline
Total            & 14,012       &  Whole                           \\ \hline
\end{tabular}}
\label{tab:dataset}
\end{table}

\subsection{Implementation Details}
\label{sec:4.2}
Our model was trained using four NVIDIA GTX 3090 Ti GPUs. We utilized the Adam optimizer~\cite{kingma2014adam} with a batch size of 8, an initial learning rate of $10^{-3}$, and a training duration of 250 epochs. 
\par In terms of pre-processing for MedLAM's training and inference, we rescaled the voxel spacing to [3,3,3] mm and standardized the cropping patch sizes to 64$\times$64$\times$64 pixels. To ensure that the scanning range was fully covered, we set the parameter $r$ as [1500, 600, 600]. 
\par In both the StructSeg HaN and WORD datasets, we start by randomly selecting 5 scans as support volumes. In the Whole-Patch Localization (WPL) approach, for each organ in these scans, we compute the six extreme coordinates and then average these coordinates and features across the 5 selected images. This creates an average representation of latent coordinates and features for each of the organ's extreme points. 
\par When employing Sub-Patch Localization (SPL), we deviate from this process. Instead of focusing on the extreme coordinates of the organ, we divide each target organ into a number of sub-patches, denoted as $i$. The value of $i$ is determined by taking the minimum span of the target organ across all support volumes and dividing it by the predefined patch distance $n$ mm. This ensures that we have a consistent number of sub-patches for each target organ across all support volumes. For each of these $i$ sub-patches, we crop, compute, and integrate the latent coordinates and features across the 5 support volumes, offering a more granular representation. As a result, we obtain the latent coordinates and features for 6$i$ points.
\par Upon obtaining the latent coordinates and features for each target point, we proceed to locate the corresponding 3D bounding boxes (bbox) in the query volume, yielding a bbox for each slice, as elaborated in Sec. \ref{sec:3.2}. To accommodate for potential variations in organ size and shape, we extend each slice's bbox prompt by [10, 10] pixels in the $x$ and $y$ directions, respectively. This allows us to ensure that the entire target structure is adequately captured for subsequent processing.
\par Upon utilizing SAM and MedSAM, the original images are subjected to a separate preprocessing routine in line with the standard procedures described in the original MedSAM methodology. This includes adjusting the slice resolution to 3 × 1024 × 1024, normalizing the intensity.
\par To guarantee the robustness of our experimental results, we repeat this random selection and averaging process five times. The final performance metrics are obtained by averaging the outcomes of these iterations.

 \subsection{Evaluation Metrics}
 To evaluate the localization accuracy of MedLAM, we apply two widely utilized metrics \citep{humpire2018efficient, xu2019efficient, lei2021contrastive}: the Intersection over Union (IoU) and Wall Distance (WD). The IoU overlap ratio is defined as:
 \begin{equation}
     IoU = \frac{\bm{V}_p \cap \bm{V}_{gt}}{\bm{V}_p \cup \bm{V}_{gt}},
 \end{equation}
where $V_p$ and $V_{gt}$ denote the volume of the predicted bbox and the ground-truth bbox, respectively. The IoU metric (0-1, 0: mismatch; 1: perfectly match) is employed to assess the relative accuracy of the predicted bounding box with respect to the ground truth. 
\par For Wall Distance (WD), it quantifies the average absolute distance difference between the boundaries of the predicted bbox and the ground-truth bbox. Specifically, it is defined as:
 \begin{equation}
     WD = \frac{1}{6}\sum{| d_{p,i} - d_{gt,i} |},
 \end{equation}
where $d_{p,i}$ and $d_{gt,i}$ are the distances of the six boundaries (top, bottom, left, right, front, back) of the predicted and ground-truth bbox, respectively. The pixel spacing is implicitly taken into account in the distance measurement.
\par As for the segmentation metrics, we use Dice Similarity Coefficient (DSC) score and 95\% Hausdorff Distance ( $\text{HD}_{95}$ ) as evaluation metrics to quantitatively evaluate our method. Dice score  (0-1, 0: mismatch; 1: perfectly match) measures the overlap of the prediction $\bm{A}$ and ground truth $\bm{B}$, and is defined as:
\begin{equation}
    DSC(\bm{A}, \bm{B})=\frac{2|\bm{A}\cap\bm{B}|}{|\bm{A}|+|\bm{B}|}\cdot 100\%
    \label{eq12}
\end{equation}
The maximum Hausdorff Distance is the maximum distance of a set $\bm{A}$ to the nearest point in the other set $\bm{B}$, defined as:
 \begin{equation}
     d_H(\bm{A},\bm{B})=\max\{ \max\limits_{a\in \bm{A}}\min_{b\in \bm{B}}d(a,b),\max\limits_{b\in \bm{B}}\min\limits_{a\in \bm{A}}d(a,b) \}
 \label{eq13}
 \end{equation}
HD$_{95}$ is similar to maximum HD, and it is based on the calculation of the 95th percentile of the distances between boundary points in $\bm{A}$ and $\bm{B}$.

\subsection{Evaluation of MedLAM}
\subsubsection{Comparison with Other Methods}
% \begin{figure*}[t!]
%   \centering
%   \includegraphics[width=0.9\textwidth]{fig/word.pdf}
%   \includegraphics[width=0.9\textwidth]{fig/structseg.pdf}
%   \caption{The mean IOU score of each organ in the WORD (top) and StructSeg (bottom) dataset.}
%   \vspace{-4mm}
%   \label{fig:iou}
% \end{figure*}
\begin{table*}[t!]
\centering
\caption{MedLAM v.s. fully-supervised \& zero-shot localization models on StructSeg Head-and-Neck and WORD datasets.}
\textcolor{black}{
\begin{subtable}{1\textwidth}
\centering
\caption{\textcolor{black}{StructSeg Head-and-Neck}}
\scalebox{1}{
\begin{tabular}{@{}ccccccc@{}}
\toprule[1.5pt]
\multirow{3}*{Organs}  & \multicolumn{3}{c}{IoU $\uparrow$ (mean$\pm$std \%)} & \multicolumn{3}{c}{WD $\downarrow$ (mean$\pm$std mm)} \\ \cmidrule(lr){2-4} \cmidrule(lr){5-7}
~ & MedLAM & nnDetection & MIU-VL & MedLAM & nnDetection & MIU-VL \\ 
~ & (5-shot) & (fully-supervised) & (zero-shot) & (5-shot) & (fully-supervised) & (zero-shot) \\ \hline
Brain Stem & 70.3$\pm$8.0 & $\bm{75.7\pm6.4}$ & 0.1$\pm$0.0 & $\bm{2.4\pm1.4}$ & 8.5$\pm$11.9 & 223.0$\pm$74.3 \\
Eye L & $\bm{66.2\pm11.5}$ & 29.0$\pm$35.8 & 0.0$\pm$0.0 & $\bm{1.7\pm1.0}$ & 7.2$\pm$13.5 & 162.8$\pm$73.9 \\
Eye R & $\bm{63.2\pm11.2}$ & 42.7$\pm$28.4 & 0.0$\pm$0.0 & $\bm{1.8\pm1.1}$ & 8.0$\pm$10.6 & 163.4$\pm$61.1 \\
Lens L & $\bm{11.5\pm11.0}$ & 10.2$\pm$16.0 & 0.0$\pm$0.0 & $\bm{1.9\pm1.1}$ & 13.8$\pm$19.5 & 220.2$\pm$89.7 \\
Lens R & 17.9$\pm$14.5 & $\bm{24.3\pm22.2}$ & 0.0$\pm$0.0 & $\bm{1.8\pm1.2}$ & 16.8$\pm$19.9 & 219.1$\pm$75.6 \\
Opt Nerve L & $\bm{29.7\pm27.0}$ & 15.4$\pm$15.6 & 0.0$\pm$0.0 & $\bm{2.0\pm1.3}$ & 9.5$\pm$14.9 & 173.0$\pm$56.3 \\
Opt Nerve R & $\bm{27.6\pm25.4}$ & 23.6$\pm$17.1 & 0.0$\pm$0.0 & $\bm{2.1\pm1.3}$ & 16.2$\pm$18.5 & 175.7$\pm$47.8 \\
Opt Chiasma & 9.7$\pm$19.4 & $\bm{32.4\pm9.4}$ & 0.5$\pm$0.0 & $\bm{2.8\pm1.5}$ & 9.0$\pm$13.0 & 177.0$\pm$25.2 \\
Temporal Lobes L & $\bm{70.9\pm8.7}$ & 45.3$\pm$37.1 & 0.5$\pm$0.1 & $\bm{3.2\pm1.8}$ & 17.0$\pm$21.0 & 163.3$\pm$33.9 \\
Temporal Lobes R & $\bm{72.8\pm10.3}$ & 26.8$\pm$33.8 & 0.0$\pm$0.2 & $\bm{3.1\pm1.8}$ & 9.5$\pm$14.8 & 189.0$\pm$64.9 \\
Pituitary & 8.9$\pm$18.5 & $\bm{25.6\pm16.8}$ & 0.3$\pm$0.0 & $\bm{2.7\pm1.6}$ & 24.4$\pm$19.5 & 229.2$\pm$84.7 \\
Parotid Gland L & $\bm{63.1\pm9.6}$ & 31.3$\pm$32.9 & 0.3$\pm$0.1 & $\bm{3.7\pm1.9}$ & 11.8$\pm$15.5 & 162.7$\pm$40.2 \\
Parotid Gland R & $\bm{62.7\pm10.3}$ & 38.7$\pm$32.0 & 0.0$\pm$0.1 & $\bm{4.0\pm2.1}$ & 13.6$\pm$20.8 & 170.1$\pm$44.6 \\
Inner Ear L & $\bm{35.6\pm8.2}$ & 34.1$\pm$25.2 & 0.0$\pm$0.0 & $\bm{2.5\pm1.0}$ & 10.3$\pm$12.2 & 182.3$\pm$34.8 \\
Inner Ear R & $\bm{33.7\pm9.4}$ & 32.1$\pm$27.7 & 0.1$\pm$0.0 & $\bm{2.7\pm1.1}$ & 15.3$\pm$16.9 & 185.4$\pm$28.0 \\
Mid Ear L & $\bm{56.5\pm12.9}$ & 31.0$\pm$32.9 & 0.1$\pm$0.1 & $\bm{3.4\pm2.0}$ & 14.0$\pm$15.7 & 193.6$\pm$72.0 \\
Mid Ear R & $\bm{56.5\pm12.8}$ & 51.1$\pm$27.5 & 0.0$\pm$0.1 & $\bm{3.5\pm2.1}$ & 15.0$\pm$17.1 & 204.6$\pm$77.5 \\
TM Joint L & $\bm{40.7\pm15.1}$ & 26.8$\pm$22.3 & 0.0$\pm$0.0 & $\bm{2.4\pm1.3}$ & 11.3$\pm$13.9 & 188.8$\pm$57.6 \\
TM Joint R & $\bm{38.7\pm15.3}$ & 26.0$\pm$18.4 & 0.5$\pm$0.0 & $\bm{2.3\pm1.2}$ & 21.4$\pm$21.4 & 200.2$\pm$64.7 \\
Spinal Cord & $\bm{57.7\pm16.7}$ & 10.4$\pm$14.8 & 0.9$\pm$0.2 & $\bm{2.7\pm1.9}$ & 11.3$\pm$15.1 & 224.5$\pm$120.1 \\
Mandible L & $\bm{80.9\pm6.2}$ & 36.1$\pm$37.3 & 1.0$\pm$0.3 & $\bm{2.5\pm1.6}$ & 13.8$\pm$17.7 & 164.2$\pm$73.1 \\
Mandible R & $\bm{80.8\pm5.9}$ & 34.3$\pm$40.0 & 1.0$\pm$0.4 & $\bm{2.6\pm1.5}$ & 21.4$\pm$20.4 & 155.2$\pm$59.6 \\ \hline
Average & $\bm{48.0\pm13.1}$ & 31.9$\pm$25.0 & 0.2$\pm$0.1 & $\bm{2.6\pm1.5}$ & 13.6$\pm$16.5 & 187.6$\pm$61.8 \\
\bottomrule[1.5pt]
\end{tabular}}
\end{subtable}}

\vspace{0.2cm}

\textcolor{black}{
\begin{subtable}{1\textwidth}
\centering
\caption{\textcolor{black}{WORD}}
\scalebox{1}{
\begin{tabular}{@{}ccccccc@{}}
\toprule[1.5pt]
\multirow{3}*{Organs}  & \multicolumn{3}{c}{IoU $\uparrow$ (mean$\pm$std \%)} & \multicolumn{3}{c}{WD $\downarrow$ (mean$\pm$std mm)} \\ \cmidrule(lr){2-4} \cmidrule(lr){5-7}
~ & MedLAM & nnDetection & MIU-VL & MedLAM & nnDetection & MIU-VL \\ 
~ & (5-shot) & (fully-supervised) & (zero-shot) & (5-shot) & (fully-supervised) & (zero-shot) \\ \hline
Liver & $\bm{73.0\pm11.6}$ & 52.7$\pm$10.2 & 10.6$\pm$2.1 & $\bm{10.5\pm9.1}$ & 23.8$\pm$31.5 & 80.4$\pm$37.2 \\
Spleen & $\bm{70.9\pm13.3}$ & 38.2$\pm$17.5 & 1.7$\pm$1.0 & $\bm{5.7\pm5.9}$ & 14.6$\pm$21.4 & 142.5$\pm$49.8 \\
Kidney L & $\bm{71.0\pm15.9}$ & 69.6$\pm$27.2 & 0.9$\pm$0.5 & $\bm{5.0\pm8.4}$ & 13.9$\pm$22.1 & 164.3$\pm$69.2 \\
Kidney R & $\bm{76.0\pm13.8}$ & 74.2$\pm$11.5 & 0.8$\pm$0.2 & $\bm{3.6\pm4.6}$ & 20.8$\pm$28.6 & 166.1$\pm$81.3 \\
Stomach & $\bm{49.1\pm14.3}$ & 36.3$\pm$23.3 & 4.7$\pm$2.1 & 17.7$\pm$12.8 & $\bm{17.0\pm25.1}$ & 102.8$\pm$25.1 \\
Gallbladder & 12.0$\pm$11.0 & $\bm{29.3\pm18.2}$ & 0.2$\pm$0.2 & $\bm{16.9\pm10.3}$ & 21.2$\pm$29.6 & 153.2$\pm$38.5 \\
Esophagus & 44.2$\pm$17.9 & $\bm{58.5\pm13.2}$ & 0.3$\pm$0.2 & $\bm{6.8\pm6.7}$ & 16.8$\pm$18.9 & 153.8$\pm$22.7 \\
Pancreas & 44.1$\pm$17.3 & $\bm{53.4\pm17.6}$ & 1.6$\pm$0.6 & $\bm{12.2\pm8.8}$ & 19.1$\pm$28.5 & 135.2$\pm$33.8 \\
Duodenum & $\bm{44.5\pm17.5}$ & 42.0$\pm$16.4 & 1.2$\pm$0.4 & $\bm{12.7\pm9.1}$ & 20.5$\pm$26.1 & 139.8$\pm$31.0 \\
Colon & $\bm{67.0\pm13.0}$ & 28.5$\pm$10.9 & 17.9$\pm$4.4 & $\bm{15.6\pm11.8}$ & 22.8$\pm$35.8 & 83.3$\pm$33.6 \\
Intestine & $\bm{62.6\pm11.0}$ & 24.6$\pm$10.7 & 12.2$\pm$4.6 & $\bm{15.4\pm9.8}$ & 20.5$\pm$29.5 & 106.0$\pm$53.1 \\
Adrenal & 40.4$\pm$17.7 & $\bm{57.3\pm12.7}$ & 0.4$\pm$0.2 & $\bm{7.6\pm5.5}$ & 16.1$\pm$26.6 & 160.2$\pm$58.0 \\
Rectum & 49.3$\pm$15.3 & $\bm{62.6\pm16.9}$ & 0.6$\pm$0.3 & $\bm{8.6\pm6.1}$ & 17.4$\pm$24.6 & 203.1$\pm$146.3 \\
Bladder & 55.4$\pm$16.9 & $\bm{63.0\pm17.0}$ & 1.2$\pm$0.7 & $\bm{8.1\pm7.6}$ & 24.1$\pm$33.2 & 167.7$\pm$96.9 \\
Head of Femur L & $\bm{76.7\pm16.7}$ & 20.0$\pm$26.8 & 1.6$\pm$2.4 & $\bm{5.3\pm16.0}$ & 14.5$\pm$23.0 & 175.1$\pm$127.4 \\
Head of Femur R & $\bm{69.4\pm14.4}$ & 47.7$\pm$34.0 & 1.5$\pm$0.8 & $\bm{5.9\pm6.9}$ & 23.8$\pm$41.5 & 163.4$\pm$117.5 \\ \hline
Average & $\bm{56.6\pm14.8}$ & 47.4$\pm$17.8 & 3.6$\pm$1.3 & $\bm{9.9\pm8.7}$ & 19.2$\pm$27.9 & 143.6$\pm$63.8 \\
\bottomrule[1.5pt]
\end{tabular}}
\end{subtable}
}
\label{tab:2}
\end{table*}

\begin{table*}[ht!]
\centering
\caption{Ablation study of Unified Anatomical Mapping (UAM) and Multi Scale Similarity (MSS) in MedLAM: IoU score (mean$\pm$std\%) on the WORD dataset}
\centering
\scalebox{0.86}{
\begin{tabular}{@{}cccccccccc@{}}
\toprule[1.5pt]
\multirow{3}*{Organs} & UAM & \XSolidBrush & UAM & UAM & UAM & UAM & UAM & UAM & UAM\\
~ & + & + & + & + & + & + & + & + & + \\
~ & \XSolidBrush & $\bm{F}^2$ + $\bm{F}^1$ + $\bm{F}^0$ & $\bm{F}^2$ & $\bm{F}^1$ & $\bm{F}^0$ & $\bm{F}^1$ + $\bm{F}^0$ & $\bm{F}^2$ + $\bm{F}^0$ & $\bm{F}^2$ + $\bm{F}^1$ & $\bm{F}^2$ + $\bm{F}^1$ + $\bm{F}^0$ \\ \hline
Liver& 57.1$\pm$10.5 & 39.7$\pm$27.4 & 76.0$\pm$8.3 & 76.7$\pm$7.9 & 71.5$\pm$10.1 & 75.0$\pm$09.1 & $\bm{78.5\pm7.7}$ & 76.6$\pm$7.7 & 73.0$\pm$11.6 \\
Spleen& 40.1$\pm$7.7 & 49.3$\pm$34.2 & 63.0$\pm$15.7 & 72.1$\pm$16.0 & 69.1$\pm$21.7 & $\bm{74.7\pm15.2}$ & 74.4$\pm$14.1 & 72.6$\pm$12.9 & 70.9$\pm$13.3 \\
Kidney L& 45.9$\pm$11.4 & 39.9$\pm$30.7 & 60.5$\pm$10.4 & 58.9$\pm$19.4 & 52.1$\pm$23.4 & 54.5$\pm$22.9 & 64.8$\pm$16.4 & 63.7$\pm$14.8 & $\bm{71.0\pm15.9}$ \\
Kidney R& 40.4$\pm$9.3 & 40.6$\pm$34.6 & 62.7$\pm$12.1 & 68.2$\pm$14.4 & 54.4$\pm$27.8 & 73.0$\pm$13.1 & $\bm{76.3\pm8.9}$ & 70.6$\pm$9.4 & 76.0$\pm$13.8 \\
Stomach& 38.7$\pm$12.2 & 25.4$\pm$21.9 & 47.4$\pm$15.7 & 43.4$\pm$15.6 & 46.4$\pm$16.7 & 47.2$\pm$17.6 & 45.9$\pm$12.9 & 47.3$\pm$14.8 & $\bm{49.1\pm14.3}$ \\
Gallbladder& 10.2$\pm$12.9 & 4.1$\pm$9.5 & 10.1$\pm$9.5 & 7.4$\pm$7.1 & 7.3$\pm$8.1 & 9.3$\pm$13.1 & 9.2$\pm$6.8 & 8.9$\pm$7.0 & $\bm{12.0\pm11.0}$ \\
Esophagus& 22.0$\pm$9.2 & 1.5$\pm$3.4 & 41.0$\pm$15.6 & 39.5$\pm$18.4 & 46.4$\pm$20.4 & 43.6$\pm$19.2 & 44.2$\pm$20.5 & 41.0$\pm$18.5 & $\bm{44.2\pm17.9}$ \\
Pancreas& 42.9$\pm$13.5 & 13.6$\pm$15.5 & 39.8$\pm$11.5 & 44.0$\pm$18.8 & 34.1$\pm$21.6 & 42.5$\pm$17.4 & $\bm{45.0\pm18.6}$ & 42.8$\pm$16.2 & 44.1$\pm$17.3 \\
Duodenum& 32.2$\pm$10.5 & 31.4$\pm$17.8 & 45.8$\pm$16.2 & 35.8$\pm$20.3 & 31.6$\pm$18.2 & 39.3$\pm$17.6 & 42.8$\pm$15.4 & $\bm{51.0\pm13.4}$ & 44.5$\pm$17.5 \\
Colon& 60.9$\pm$13.4 & 42.8$\pm$25.6 & 63.7$\pm$12.9 & 65.3$\pm$12.1 & 61.0$\pm$11.1 & 65.2$\pm$11.8 & 64.5$\pm$11.7 & 63.2$\pm$12.5 & $\bm{67.0\pm13.0}$ \\
Intestine& 58.7$\pm$11.6 & 27.7$\pm$21.1 & 63.1$\pm$13.0 & 61.8$\pm$16.8 & 62.0$\pm$17.0 & 61.9$\pm$17.6 & $\bm{65.2\pm13.1}$ & 63.9$\pm$14.0 & 62.6$\pm$11.0 \\
Adrenal& 29.8$\pm$17.0 & 10.5$\pm$15.6 & 36.8$\pm$16.1 & 30.5$\pm$18.8 & 32.2$\pm$19.7 & 35.0$\pm$19.7 & 38.1$\pm$15.7 & 38.6$\pm$13.5 & $\bm{40.4\pm17.7}$ \\
Rectum& 34.6$\pm$10.5 & 47.9$\pm$18.2 & 49.1$\pm$12.4 & 47.5$\pm$11.6 & 32.1$\pm$19.4 & 47.9$\pm$12.7 & 51.8$\pm$12.7 & $\bm{53.3\pm13.7}$ & 49.3$\pm$15.3 \\
Bladder& 33.9$\pm$15.2 & 45.8$\pm$23.0 & 49.5$\pm$18.1 & 55.3$\pm$21.1 & 49.6$\pm$20.1 & 52.3$\pm$22.2 & 50.8$\pm$20.4 & $\bm{55.4\pm16.9}$ & 55.4$\pm$16.9 \\
Head of Femur L& 43.7$\pm$10.8 & 71.9$\pm$22.4 & 71.8$\pm$8.5 & 80.2$\pm$8.7 & $\bm{84.1\pm8.9}$ & 83.1$\pm$8.9 & 81.6$\pm$9.1 & 77.9$\pm$8.5 & 76.7$\pm$16.7 \\
Head of Femur R& 36.2$\pm$10.9 & 55.3$\pm$29.4 & 63.8$\pm$13.3 & 70.9$\pm$12.5 & $\bm{75.9\pm13.7}$ & 74.0$\pm$14.0 & 70.2$\pm$14.2 & 68.3$\pm$13.9 & 69.4$\pm$14.4 \\
\hline
Average& 39.2$\pm$11.7 & 34.2$\pm$21.9 & 52.8$\pm$13.1 & 53.6$\pm$15.0 & 50.6$\pm$17.4 & 54.9$\pm$15.8 & 56.5$\pm$13.6 & 55.9$\pm$13.2 & $\bm{56.6\pm14.8}$ \\
\bottomrule[1.5pt]
\end{tabular}}
\label{tab:3}
\end{table*}

In this section, with the purpose of validating the performance of our universal localization model MedLAM. We have conducted a comparative analysis between  MedLAM and two well-established methodologies in the medical image detection field: 1) nnDetection~\cite{baumgartner2021nndetection}, which serves as a highly robust fully-supervised medical image detection baseline; 2) MIU-VL~\cite{qin2023medical}, a pioneer work utilizing visual-language foundation model for zero-shot medical image localization.

For nnDetection, we followed the automated data processing strategy and performed five-fold cross-training approach to identify the optimal network hyperparameters. Regarding MIU-VL, we follow the procedures outlined in the original paper, using GLIP \citep{li2022grounded} as the foundation model for zero-shot detection.  In order to generate the required text prompts for MIU-VL, we utilized ChatGPT \citep{ouyang2022training} to generate prompts containing descriptions of [COLOR], [SHAPE], [LOCATION], and [CLASS] for different organs.  Since the MIU-VL network is designed to directly handle 2D images, we performed the detection process on 3D medical images slice by slice. Subsequently, we reconstructed the obtained bounding boxes back into the 3D image space.

The final experimental results are presented in Table~\ref{tab:2}. MedLAM significantly outperforms both nnDetection and MIU-VL on the StructSeg HaN and WORD datasets. Specifically, on the StructSeg HaN dataset, MedLAM achieved an average IoU of 48\%, substantially surpassing nnDetection's 16.1\%. We also observed that MIU-VL was only able to delineate the entire body area regardless of the organ-specific prompt, highlighting a substantial gap between foundation models pre-trained on natural images and their applicability in the medical domain. These results also indicate that such natural image foundation models may struggle to provide meaningful localization in low-contrast medical images. In contrast, our MedLAM model, which incorporates strong medical domain knowledge, is capable of localizing any anatomical structure with just a few annotated examples, exceeding the performance of fully supervised models.

To be more specific, in the StructSeg HaN dataset, the organ with the largest average WD (Wall Distance) is the right parotid gland, which does not exceed 4mm. This suggests that MedLAM can perform highly accurate localization for head and neck organs. On the other hand, for the WORD dataset, the abdominal soft tissue organs have blurry boundaries and significant deformations, resulting in larger average WD values. However, the WD for each organ is generally within 15mm.

In summary, MedLAM exhibits robust localization performance and proves to be a versatile foundational model for localizing any type of structure in 3D medical imaging. For smaller organs, although the IoU scores are relatively lower, the subsequent preprocessing step of expanding the 3D bounding box ensures that these organs are adequately captured for segmentation. This makes the model practically applicable across a broad spectrum of organ sizes. This approach provides a pragmatic solution to the challenge of localizing smaller organs with MedLAM, thus achieving a balanced performance across different organ sizes.

\begin{table*}[t!]
    \centering
    \caption{Comparison of segmentation performance between MedLAM-generated bounding box prompts and manually annotated prompts on the StructSeg Head-and-Neck dataset. Metrics include Dice Similarity Coefficient (DSC) and Hausdorff Distance (HD$_{95}$).}
    \textcolor{black}{
    \begin{subtable}{1\textwidth}
    \centering
    \caption{\textcolor{black}{StructSeg Head-and-Neck}}
    \scalebox{0.95}{
    \begin{tabular}{ccccccccc}
    \toprule[1.5pt]
    \multirow{3}*{Organs} & \multicolumn{4}{c}{DSC $\uparrow$ (mean$\pm$std \%)} & \multicolumn{4}{c}{HD$_{95}$ $\downarrow$ (mean$\pm$std mm)} \\ \cmidrule(lr){2-5} \cmidrule(lr){6-9}
    ~ & \multicolumn{2}{c}{MedLAM} & \multicolumn{2}{c}{Manual Prompt} & \multicolumn{2}{c}{MedLAM} & \multicolumn{2}{c}{Manual Prompt} \\ \cmidrule(lr){2-3} \cmidrule(lr){4-5} \cmidrule(lr){6-7} \cmidrule(lr){8-9}
    ~ & SAM & MedSAM & SAM & MedSAM & SAM & MedSAM & SAM & MedSAM\\\hline
    Brain Stem       & 63.5$\pm$6.3       & $\bm{73.3\pm5.0}$  & 66.2$\pm$2.7       & $\bm{75.9\pm2.7}$ & 8.4$\pm$1.6       & $\bm{6.2\pm1.2}$  & 8.6$\pm$1.0      & $\bm{6.1\pm1.1}$ \\
    Eye L            & 64.5$\pm$6.5       & $\bm{69.3\pm5.9}$  & 65.7$\pm$5.6       & $\bm{70.0\pm4.6}$ & 8.0$\pm$1.5       & $\bm{5.8\pm1.0}$  & 7.9$\pm$1.3      & $\bm{5.9\pm1.1}$ \\
    Eye R            & 67.3$\pm$5.8       & $\bm{69.4\pm5.5}$  & 68.9$\pm$4.8       & $\bm{69.1\pm5.2}$ & 7.6$\pm$1.9       & $\bm{5.7\pm1.2}$  & 7.4$\pm$1.2      & $\bm{6.2\pm1.1}$ \\
    Lens L           & 15.9$\pm$7.8       & $\bm{16.0\pm6.7}$  & $\bm{22.7\pm5.6}$  & 19.2$\pm$3.5      & $\bm{8.0\pm1.5}$  & 8.6$\pm$1.7       & 10.0$\pm$2.3     & $\bm{9.5\pm1.3}$ \\
    Lens R           & 13.8$\pm$8.8       & $\bm{14.0\pm5.5}$  & $\bm{22.8\pm7.7}$  & 17.3$\pm$4.6      & $\bm{9.3\pm1.5}$  & 9.6$\pm$1.3       & $\bm{9.7\pm2.1}$ & 10.9$\pm$2.3     \\
    Opt Nerve L      & $\bm{23.7\pm6.1}$  & 23.5$\pm$6.2       & 32.6$\pm$9.3       & $\bm{34.7\pm5.4}$ & 9.0$\pm$1.7       & $\bm{8.6\pm1.1}$  & 10.1$\pm$2.5     & $\bm{8.4\pm1.3}$ \\
    Opt Nerve R      & $\bm{27.8\pm10.2}$ & 26.3$\pm$6.6       & 28.5$\pm$6.7       & $\bm{32.4\pm5.2}$ & 9.2$\pm$3.2       & $\bm{9.0\pm1.4}$  & 11.3$\pm$2.2     & $\bm{9.4\pm1.0}$ \\
    Opt Chiasma      & 11.4$\pm$10.6      & $\bm{14.4\pm11.4}$ & 39.8$\pm$10.2      & $\bm{39.8\pm8.1}$ & 10.2$\pm$1.4      & $\bm{8.8\pm1.6}$  & 7.2$\pm$1.7      & $\bm{6.6\pm1.2}$ \\
    Temporal Lobes L & 28.2$\pm$15.2      & $\bm{78.3\pm3.5}$  & 36.8$\pm$16.6      & $\bm{83.5\pm1.9}$ & 16.4$\pm$3.5      & $\bm{9.0\pm2.1}$  & 13.7$\pm$3.5     & $\bm{6.6\pm1.4}$ \\
    Temporal Lobes R & 24.1$\pm$17.4      & $\bm{78.0\pm4.3}$  & 30.7$\pm$17.8      & $\bm{84.4\pm1.4}$ & 18.2$\pm$4.2      & $\bm{10.1\pm2.4}$ & 15.1$\pm$4.1     & $\bm{6.1\pm2.5}$ \\
    Pituitary        & $\bm{12.5\pm10.7}$ & 10.2$\pm$9.2       & $\bm{36.6\pm16.2}$ & 27.5$\pm$12.5     & $\bm{8.7\pm2.2}$  & 8.8$\pm$1.8       & $\bm{7.4\pm2.0}$ & 10.0$\pm$3.6     \\
    Parotid Gland L  & 15.5$\pm$11.9      & $\bm{59.6\pm6.5}$  & 29.5$\pm$7.7       & $\bm{64.4\pm3.4}$ & 17.6$\pm$3.0      & $\bm{10.3\pm1.9}$ & 16.0$\pm$1.9     & $\bm{9.3\pm1.9}$ \\
    Parotid Gland R  & 17.2$\pm$9.9       & $\bm{57.1\pm6.8}$  & 29.6$\pm$6.5       & $\bm{64.9\pm5.0}$ & 19.8$\pm$5.0      & $\bm{13.5\pm3.5}$ & 14.0$\pm$2.1     & $\bm{9.5\pm1.6}$ \\
    Inner Ear L      & 40.4$\pm$11.8      & $\bm{42.3\pm9.9}$  & 59.7$\pm$11.7      & $\bm{62.1\pm9.7}$ & 9.4$\pm$1.7       & $\bm{7.6\pm1.5}$  & 7.6$\pm$2.1      & $\bm{5.9\pm1.8}$ \\
    Inner Ear R      & $\bm{48.9\pm9.5}$  & 45.9$\pm$11.2      & $\bm{67.3\pm10.4}$ & 66.9$\pm$6.1      & 8.4$\pm$1.7       & $\bm{7.0\pm1.7}$  & 6.8$\pm$1.9      & $\bm{5.4\pm1.2}$ \\
    Mid Ear L        & $\bm{64.6\pm14.3}$ & 59.7$\pm$9.6       & $\bm{71.0\pm11.7}$ & 67.5$\pm$6.3      & 7.9$\pm$2.6       & $\bm{7.7\pm1.6}$  & 6.4$\pm$2.8      & $\bm{6.1\pm1.1}$ \\
    Mid Ear R        & $\bm{64.7\pm13.1}$ & 59.3$\pm$11.2      & $\bm{73.3\pm7.7}$  & 65.0$\pm$8.9      & 8.4$\pm$2.4       & $\bm{7.6\pm1.4}$  & $\bm{6.2\pm1.9}$ & 6.6$\pm$1.2      \\
    TM Joint L       & 38.3$\pm$10.1      & $\bm{39.0\pm10.9}$ & 59.7$\pm$12.8      & $\bm{61.2\pm4.9}$ & 10.2$\pm$2.1      & $\bm{7.7\pm1.6}$  & 6.8$\pm$1.9      & $\bm{6.0\pm0.7}$ \\
    TM Joint R       & $\bm{41.5\pm10.0}$ & 38.3$\pm$9.5       & 59.2$\pm$17.8      & $\bm{60.0\pm7.7}$ & 9.6$\pm$1.6       & $\bm{8.1\pm1.4}$  & 8.2$\pm$1.9      & $\bm{6.3\pm1.1}$ \\
    Spinal Cord      & 27.9$\pm$8.3       & $\bm{34.7\pm6.9}$  & 38.0$\pm$7.7       & $\bm{40.3\pm6.1}$ & 15.5$\pm$2.7      & $\bm{13.3\pm2.9}$ & 10.7$\pm$1.3     & $\bm{8.8\pm0.9}$ \\
    Mandible L       & $\bm{78.0\pm4.9}$  & 66.7$\pm$6.1       & $\bm{86.0\pm2.1}$  & 76.8$\pm$4.7      & $\bm{11.1\pm2.8}$ & 11.8$\pm2.1$      & $\bm{5.5\pm1.7}$ & 9.0$\pm$4.5      \\
    Mandible R       & $\bm{71.4\pm4.0}$  & 66.0$\pm$4.7       & $\bm{81.2\pm2.7}$  & 75.0$\pm$4.5      & 22.6$\pm$4.7      & $\bm{12.3\pm2.6}$ & 21.6$\pm$5.4     & $\bm{9.6\pm3.5}$ \\ \hline
    Average          & 39.1$\pm$9.7       & $\bm{47.3\pm7.4}$ & 50.3$\pm$9.2       & $\bm{57.2\pm5.6}$ & 11.5$\pm$2.5      & $\bm{9.0\pm1.8}$  & 9.9$\pm$2.2      & $\bm{7.6\pm1.7}$ \\ 
    \bottomrule[1.5pt]
    \end{tabular}}
    \end{subtable}}
    
    \vspace{0.2cm}
    
    \textcolor{black}{
    \begin{subtable}{1\textwidth}
    \centering
    \caption{\textcolor{black}{WORD}}
    \scalebox{0.95}{
    \begin{tabular}{ccccccccc} \bottomrule[1.5pt]
    \multirow{3}*{Organs} & \multicolumn{4}{c}{DSC $\uparrow$ (mean$\pm$std \%)} & \multicolumn{4}{c}{HD$_{9}$ $\downarrow$ (mean$\pm$std mm)} \\ \cmidrule(lr){2-5} \cmidrule(lr){6-9}
    ~ & \multicolumn{2}{c}{MedLAM} & \multicolumn{2}{c}{Manual Prompt} & \multicolumn{2}{c}{MedLAM} & \multicolumn{2}{c}{Manual Prompt} \\ \cmidrule(lr){2-3} \cmidrule(lr){4-5} \cmidrule(lr){6-7} \cmidrule(lr){8-9}
    ~ & SAM & MedSAM & SAM & MedSAM & SAM & MedSAM & SAM & MedSAM\\\hline
    Liver           & $\bm{66.0\pm10.1}$ & 23.8$\pm$8.6      & $\bm{84.2\pm6.3}$  & 46.6$\pm$14.7     & 53.6$\pm$10.4       & $\bm{51.5\pm8.6}$  & $\bm{21.8\pm11.9}$ & 36.8$\pm$11.6      \\
    Spleen          & $\bm{61.7\pm14.4}$ & 36.3$\pm$13.2     & $\bm{85.3\pm5.1}$  & 65.0$\pm$6.3      & $\bm{30.1\pm16.1}$  & 30.5$\pm$14.4      & $\bm{10.8\pm6.5}$  & 16.2$\pm$4.9       \\
    Kidney L        & $\bm{82.1\pm16.0}$ & 70.7$\pm$18.8     & $\bm{92.1\pm1.6}$  & 84.1$\pm$5.2      & 19.2$\pm$18.6       & $\bm{18.5\pm16.9}$ & $\bm{6.1\pm1.9}$   & 8.2$\pm$1.6        \\
    Kidney R        & $\bm{88.3\pm4.8}$  & 77.3$\pm$6.2      & $\bm{92.9\pm1.6}$  & 86.4$\pm$3.1      & $\bm{12.7\pm5.7}$   & 13.4$\pm$4.1       & $\bm{6.0\pm1.6}$   & 7.7$\pm$1.2        \\
    Stomach         & $\bm{44.6\pm15.0}$ & 37.2$\pm$14.7     & 77.1$\pm$10.6      & $\bm{80.3\pm4.7}$ & $\bm{50.1\pm27.0}$  & 51.7$\pm$27.6      & 12.7$\pm$4.2       & $\bm{11.5\pm3.6}$  \\
    Gallbladder     & $\bm{13.1\pm17.3}$ & 10.2$\pm$13.4     & $\bm{72.7\pm11.1}$ & 68.8$\pm$7.8      & 36.8$\pm$19.8       & $\bm{36.7\pm20.4}$ & $\bm{6.0\pm1.8}$   & 6.4$\pm$1.6        \\
    Esophagus       & $\bm{36.6\pm14.7}$ & 27.8$\pm$12.9     & $\bm{67.0\pm6.4}$  & 63.1$\pm$7.7      & 21.1$\pm$12.7       & $\bm{19.8\pm11.6}$ & 6.9$\pm$4.0        & $\bm{6.7\pm3.2}$   \\
    Pancreas        & $\bm{29.7\pm14.1}$ & 21.4$\pm$9.7      & $\bm{64.4\pm7.7}$  & 46.9$\pm$11.6     & $\bm{32.6\pm13.5}$  & 33.1$\pm$10.8      & $\bm{15.9\pm6.4}$  & 18.3$\pm$5.5       \\
    Duodenum        & $\bm{26.0\pm11.4}$ & 21.1$\pm$7.0      & $\bm{54.1\pm13.2}$ & 51.0$\pm$11.9     & 31.1$\pm$9.7        & $\bm{30.9\pm10.4}$ & $\bm{16.0\pm4.6}$  & 17.1$\pm$5.9       \\
    Colon           & 25.6$\pm$9.6       & $\bm{26.6\pm8.9}$ & 41.8$\pm$6.8       & $\bm{44.1\pm8.8}$ & $\bm{57.8\pm17.8}$  & 60.6$\pm$17.8      & 49.4$\pm$15.3      & $\bm{49.1\pm13.9}$ \\
    Intestine       & $\bm{37.5\pm7.6}$  & 34.1$\pm$8.4      & $\bm{61.4\pm6.9}$  & 52.5$\pm$8.0      & 47.1$\pm$10.5       & $\bm{44.1\pm10.0}$ & 21.7$\pm$4.2       & $\bm{21.0\pm4.0}$  \\
    Adrenal         & 3.3$\pm$3.9        & $\bm{10.0\pm6.2}$ & 17.4$\pm$8.8       & $\bm{26.5\pm5.9}$ & 30.8$\pm$8.7        & $\bm{29.9\pm7.5}$  & 23.4$\pm$4.3       & $\bm{22.6\pm2.7}$  \\
    Rectum          & $\bm{50.1\pm17.9}$ & 46.0$\pm$18.7     & 75.5$\pm$4.0       & $\bm{80.0\pm3.6}$ & 18.6$\pm$7.0        & $\bm{17.4\pm7.2}$  & 7.6$\pm$0.9        & $\bm{4.8\pm0.6}$   \\
    Bladder         & $\bm{65.3\pm26.6}$ & 59.1$\pm$23.1     & $\bm{83.0\pm15.5}$ & 82.9$\pm$8.0      & 20.5$\pm$10.2       & $\bm{19.8\pm8.8}$  & 7.5$\pm$3.5        & $\bm{6.9\pm2.0}$   \\
    Head of Femur L & $\bm{81.7\pm3.9}$  & 71.5$\pm$4.1      & $\bm{90.5\pm2.7}$  & 80.3$\pm$2.8      & 16.6$\pm$3.3        & $\bm{14.9\pm1.6}$  & $\bm{9.2\pm4.7}$   & 12.3$\pm$2.0       \\
    Head of Femur R & $\bm{80.1\pm3.3}$  & 74.3$\pm$4.4      & $\bm{89.1\pm3.4}$  & 83.2$\pm$2.7      & 19.5$\pm$2.6        & $\bm{13.4\pm1.0}$  & 12.6$\pm$6.1       & $\bm{12.3\pm7.7}$  \\ \hline
    Average         & $\bm{49.5\pm11.9}$ & 40.5$\pm$11.1     & $\bm{71.8\pm7.0}$  & 65.1$\pm$7.0      & 31.1$\pm$12.1       & $\bm{30.4\pm11.2}$ & $\bm{14.6\pm5.1}$  & 16.1$\pm$4.5       \\
    \bottomrule[1.5pt]
    \end{tabular}}
    \end{subtable}
    }
    \label{tab:4}
    \end{table*}

\begin{table*}[t!]
\centering
\caption{Comparison of DSC scores $\downarrow$ (mean $\pm$ std\%) for MedLSAM using different localization strategies on the StructSeg Head-and-Neck dataset: Whole-Patch Localization (WPL) and Sub-Patch Localization (SPL) with varying slice intervals $n$.}
\centering
\scalebox{0.95}{
\begin{tabular}{ccccccccc}
\toprule[1.5pt]
\multirow{3}*{Organs} & \multicolumn{4}{c}{SAM} & \multicolumn{4}{c}{MedSAM} \\ \cmidrule(lr){2-5} \cmidrule(lr){6-9}
~ & \multicolumn{3}{c}{SPL} & WPL & \multicolumn{3}{c}{SPL} & WPL  \\ \cmidrule(lr){2-4} \cmidrule(lr){5-5} \cmidrule(lr){6-8} \cmidrule(lr){9-9}
~ & $n=6mm$ & $n=15mm$ & $n=30mm$ & \XSolidBrush & $n=6mm$ & $n=15mm$ & $n=30mm$ & \XSolidBrush  \\ \hline
Brain Stem & $\bm{63.5\pm6.3}$& 60.1$\pm$6.7&51.5$\pm$6.0& 50.8$\pm$6.2&$\bm{73.3\pm5.0}$&  71.9$\pm$4.8&   67.4$\pm$4.2&  66.6$\pm$4.4\\
Eye L &$\bm{64.5\pm6.5}$& 61.2$\pm$6.7&61.2$\pm$6.7& 61.2$\pm$6.7&  $\bm{69.3\pm5.9}$&  66.4$\pm$5.7&   66.4$\pm$5.7&  66.4$\pm$5.7\\
Eye R &$\bm{67.3\pm5.8}$& 64.5$\pm$6.6&64.5$\pm$6.6& 64.5$\pm$6.6&  $\bm{69.4\pm5.5}$&  67.8$\pm$6.0&   67.8$\pm$6.0&  67.8$\pm$6.0\\
Lens L &$\bm{15.9\pm7.8}$& 15.7$\pm$7.7&15.7$\pm$7.7& 15.7$\pm$7.7&  $\bm{16.0\pm6.7}$&  15.8$\pm$6.7&   15.8$\pm$6.7&  15.8$\pm$6.7\\
Lens R & 13.8$\pm$8.8&$\bm{14.1\pm8.6}$&13.5$\pm$8.7& 13.5$\pm$8.7&  $\bm{14.0\pm5.5}$&  $\bm{14.0\pm5.6}$&   13.7$\pm$5.5&  13.7$\pm$5.5\\
Opt Nerve L & 23.7$\pm$6.1&$\bm{23.9\pm5.9}$&23.7$\pm$6.1& 23.7$\pm$6.1&  23.5$\pm$6.2&  $\bm{23.6\pm6.2}$&   23.5$\pm$6.2&  23.5$\pm$6.2\\
Opt Nerve R &$\bm{27.8\pm10.2}$&$\bm{27.8\pm10.2}$&$\bm{27.8\pm10.2}$&$\bm{27.8\pm10.2}$&  $\bm{26.3\pm6.6}$&  $\bm{26.3\pm6.6}$&   $\bm{26.3\pm6.6}$&  $\bm{26.3\pm6.6}$\\
Opt Chiasma &$\bm{11.4\pm10.6}$&$\bm{11.4\pm10.6}$& $\bm{11.4\pm10.6}$& $\bm{11.4\pm10.6}$& $\bm{14.4\pm11.4}$& $\bm{14.4\pm11.4}$&  $\bm{14.4\pm11.4}$& $\bm{14.4\pm11.4}$\\
Temporal Lobes L &$\bm{28.2\pm15.2}$&21.7$\pm$14.8& 25.7$\pm$13.4& 25.7$\pm$13.4&  $\bm{78.3\pm3.5}$&  73.1$\pm$3.5&   71.0$\pm$3.6&  71.0$\pm$3.6\\
Temporal Lobes R &$\bm{24.1\pm17.4}$&18.2$\pm$17.7& 18.5$\pm$16.4& 18.5$\pm$16.4&  $\bm{78.0\pm4.3}$&  75.9$\pm$3.1&   71.0$\pm$3.5&  71.0$\pm$3.5\\
Pituitary &12.5$\pm$10.7&$\bm{12.6\pm10.7}$&$\bm{12.6\pm10.7}$&$\bm{12.6\pm10.7}$&  10.2$\pm$9.2&  $\bm{10.3\pm9.4}$&   $\bm{10.3\pm9.4}$&  $\bm{10.3\pm9.4}$\\
Parotid Gland L &$\bm{15.5\pm11.9}$& 10.4$\pm$9.6& 5.5$\pm$6.6& 5.5$\pm$6.6&  $\bm{59.6\pm6.5}$&  56.4$\pm$6.6&   48.1$\pm$7.3&  48.1$\pm$7.3\\
Parotid Gland R &$\bm{17.2\pm9.9}$&13.1$\pm$10.5& 7.7$\pm$7.7& 7.7$\pm$7.7&  $\bm{57.1\pm6.8}$ &  54.6$\pm$6.8&   45.3$\pm$7.8&  45.3$\pm$7.8\\
Inner Ear L &$\bm{40.4\pm11.8}$&$\bm{40.4\pm11.8}$&$\bm{40.4\pm11.8}$&$\bm{40.4\pm11.8}$&  $\bm{42.3\pm9.9}$&  $\bm{42.3\pm9.9}$&  $\bm{42.3\pm9.9}$ & $\bm{42.3\pm9.9}$ \\
Inner Ear R &$\bm{48.9\pm9.5}$&$\bm{48.9\pm9.5}$&$\bm{48.9\pm9.5}$&$\bm{48.9\pm9.5}$& $\bm{45.9\pm11.2}$ & $\bm{45.9\pm11.2}$ & $\bm{45.9\pm11.2}$ & $\bm{45.9\pm11.2}$ \\
Mid Ear L &$\bm{64.6\pm14.3}$&62.3$\pm$14.7& 62.3$\pm$14.7& 62.3$\pm$14.7& $\bm{59.7\pm9.6}$ & 55.8$\pm$10.0&  55.7$\pm$10.0& 55.7$\pm$10.0\\
Mid Ear R &$\bm{64.7\pm13.1}$&62.4$\pm$13.1& 62.4$\pm$13.0& 62.4$\pm$13.0& $\bm{59.3\pm11.2}$ & 53.6$\pm$12.1&  53.4$\pm$11.9& 53.4$\pm$11.9\\
TM Joint L &38.3$\pm$10.1&38.3$\pm$10.1& 38.3$\pm$10.1& 38.3$\pm$10.1& $\bm{39.0\pm10.9}$ & $\bm{39.0\pm10.9}$ & $\bm{39.0\pm10.9}$ & $\bm{39.0\pm10.9}$ \\
TM Joint R &41.5$\pm$10.0&41.5$\pm$10.0& 41.5$\pm$10.0& 41.5$\pm$10.0&  $\bm{38.3\pm9.5}$&  $\bm{38.3\pm9.5}$&  $\bm{38.3\pm9.5}$& $\bm{38.3\pm9.5}$\\
Spinal Cord & 27.9$\pm$8.3& 27.7$\pm$8.2&27.8$\pm$8.3& 7.8$\pm$2.9&  $\bm{34.7\pm6.9}$&  $\bm{34.7\pm6.9}$& $\bm{34.7\pm6.9}$& 11.7$\pm$3.9\\
Mandible L & $\bm{78.0\pm4.9}$& $\bm{78.0\pm5.2}$&68.0$\pm$6.2& 47.0$\pm$5.3& $\bm{66.7\pm6.1}$ &  $\bm{66.7\pm5.2}$ &  46.7$\pm$5.8&  22.9$\pm$4.1\\
Mandible R & $\bm{71.4\pm4.0}$ & 69.0$\pm$4.4&58.5$\pm$3.9& 40.2$\pm$4.6& $\bm{66.0\pm4.7}$ &  64.7$\pm$4.9&   48.2$\pm$5.0&  25.9$\pm$3.2\\
\hline
Average & $\bm{39.1\pm9.7}$ & 37.4$\pm$9.7 & 35.8$\pm$9.2 & 33.1$\pm$9.1 & $\bm{47.3\pm7.4}$ & 46.0$\pm$7.4 & 43.0$\pm$7.5 & 39.8$\pm$7.2 \\ 
\bottomrule[1.5pt]
\end{tabular}}
\label{tab:5}
\end{table*}

\begin{figure*}[ht!]
    \centering
    \includegraphics[width=1\linewidth]{seg_result.pdf}
    \caption{Visualization examples of segmentation results on WORD and StructSeg Head-and-Neck datasets using pre-trained MedSAM and SAM, post landmark localization with MedLAM.}
    \label{fig:seg_result}
\end{figure*}

\subsubsection{Ablation Study of MedLAM}
To assess the contributions of individual components within MedLAM, we conduct a series of experiments on the WORD dataset. Specifically, we evaluate MedLAM under the following conditions: 1) using only Unified Anatomical Mapping (UAM) for coarse localization; 2) relying solely on Multi Scale Similarity (MSS) for pinpointing the most similar point across the entire query image; 3) employing RDR for initial coarse localization, followed by MSS with and without features from each scale. The IoU scores obtained from these experiments are presented in Table \ref{tab:3}.
As can be observed, when using either RDR or MSS independently, the average IoU is below 50\%, standing at 39.2\% and 34.2\%, respectively. This could be attributed to the following reasons: 1) RDR operates under the assumption of anatomical consistency, which may not hold true if the query volume undergoes significant deformation; 2) MSS focuses on the similarity of local features but lacks an understanding of global anatomical structures, making it susceptible to interference from other areas in the image that are similar to the query point. This results in the identification of false positives, as evidenced by the high variance in the MSS detection results.

However, when initiating with a coarse localization point identified by RDR and subsequently refining it with MSS, the detection accuracy is significantly improved. As shown in Table \ref{tab:3}, the average IoU can be increased from 39.2\% to 56.6\% when using multi-scale features for MSS-based refining. Furthermore, using individual feature maps for $\bm{F}^i$ refining also improves the average IoU to over 50\%, validating the effectiveness of combining global anatomical features with local pixel features.
\subsection{Evaluation of MedLSAM}
\subsubsection{Effectness of MedLSAM}

% \begin{table*}[t!]
% \centering
% \caption{Ablation study in DSC score (\%) and HD$_{95}$ (mm) on $\tau$ values in MSD (multi similarity denoising).}
% \textcolor{black}{
% \begin{subtable}{1\textwidth}
% \centering
% \caption{\textcolor{black}{StructSeg Head-and-Neck}}
% \scalebox{0.95}{
% \begin{tabular}{ccccccccc}
% \toprule[1.5pt]
% \multirow{3}*{Organs} & \multicolumn{4}{c}{DSC $\uparrow$ (mean$\pm$std \%)} & \multicolumn{4}{c}{HD$_{95}$ $\downarrow$ (mean$\pm$std mm)} \\ \cmidrule(lr){2-5} \cmidrule(lr){6-9}
% ~ & \multicolumn{2}{c}{MedLAM} & \multicolumn{2}{c}{Manual Prompt} & \multicolumn{2}{c}{MedLAM} & \multicolumn{2}{c}{Manual Prompt} \\ \cmidrule(lr){2-3} \cmidrule(lr){4-5} \cmidrule(lr){6-7} \cmidrule(lr){8-9}
% ~ & SAM & MedSAM & SAM & MedSAM & SAM & MedSAM & SAM & MedSAM\\\hline
% Brain Stem & 60.1$\pm$6.7& $\bm{71.9\pm4.8}$ & 66.2$\pm$2.7& $\bm{75.9\pm2.7}$ & $\bm{9.5\pm1.8}$ & 6.5$\pm$1.2& 8.6$\pm$1.0& $\bm{6.1\pm1.1}$\\
% Eye L & 61.2$\pm$6.7& $\bm{66.4\pm5.7}$ & 65.7$\pm$5.6& $\bm{70.0\pm4.6}$ & 9.1$\pm$1.7& $\bm{6.3\pm1.0}$ & 7.9$\pm$1.3& $\bm{5.9\pm1.1}$\\
% Eye R & 64.5$\pm$6.6& $\bm{67.8\pm6.0}$ & 68.9$\pm$4.8& $\bm{69.1\pm5.2}$ & 8.0$\pm$1.9& $\bm{6.0\pm1.0}$ & 7.4$\pm$1.2& $\bm{6.2\pm1.1}$\\
% Lens L & 15.7$\pm$7.7 & $\bm{15.8\pm6.7}$ & $\bm{22.7\pm5.6}$ & 19.2$\pm$3.5& $\bm{8.0\pm1.4}$ & 8.5$\pm$1.6 & $\bm{10.0\pm2.3}$ & 9.5$\pm$1.3\\
% Lens R & $\bm{14.1\pm8.6}$ & 14.0$\pm$5.6& $\bm{22.8\pm7.7}$ & 17.3$\pm$4.6& $\bm{9.4\pm1.6}$& 9.5$\pm$1.3 & $\bm{9.7\pm2.1}$ & 10.9$\pm$2.3\\
% Opt Nerve L & $\bm{23.9\pm5.9}$ & 23.6$\pm$6.2 & 32.6$\pm$9.3& $\bm{34.7\pm5.4}$ & 8.7$\pm$1.4 & $\bm{8.5\pm0.9}$ & 10.1$\pm$2.5& $\bm{8.4\pm1.3}$\\
% Opt Nerve R & $\bm{27.8\pm10.2}$ & 26.3$\pm$6.6 & 28.5$\pm$6.7& $\bm{32.4\pm5.2}$ & 9.2$\pm$3.1 & $\bm{9.0\pm1.4}$ & 11.3$\pm$2.2& $\bm{9.4\pm1.0}$\\
% Opt Chiasma & 11.4$\pm$10.6& $\bm{14.4\pm11.4}$ & 39.8$\pm$10.2 & $\bm{39.8\pm8.1}$& 10.2$\pm$1.3& $\bm{8.8\pm1.5}$ & 7.2$\pm$1.7 & $\bm{6.6\pm1.2}$\\
% Temporal Lobes L & 21.7$\pm$14.8& $\bm{73.1\pm3.5}$ & 36.8$\pm$16.6& $\bm{83.5\pm1.9}$ & 17.0$\pm$3.1& $\bm{12.0\pm1.4}$ & 13.7$\pm$3.5& $\bm{6.6\pm1.4}$\\
% Temporal Lobes R & 18.2$\pm$17.7& $\bm{75.9\pm3.1}$ & 30.7$\pm$17.8& $\bm{84.4\pm1.4}$ & 19.4$\pm$4.8& $\bm{10.5\pm2.0}$ & 15.1$\pm$4.1& $\bm{6.1\pm2.5}$\\
% Pituitary & $\bm{12.6\pm10.7}$ & 10.3$\pm$9.4& $\bm{36.6\pm16.2}$ & 27.5$\pm$12.5& 8.7$\pm$2.1& $\bm{8.7\pm1.8}$ & $\bm{7.4\pm2.0}$ & 10.0$\pm$3.6\\
% Parotid Gland L & 10.4$\pm$9.6& $\bm{56.4\pm6.6}$ & 29.5$\pm$7.7& $\bm{64.4\pm3.4}$ & 19.7$\pm$2.7& $\bm{11.4\pm2.2}$ & 16.0$\pm$1.9& $\bm{9.3\pm1.9}$\\
% Parotid Gland R & 13.1$\pm$10.5& $\bm{54.6\pm6.8}$ & 29.6$\pm$6.5& $\bm{64.9\pm5.0}$ & 21.4$\pm$5.2& $\bm{14.7\pm3.6}$ & 14.0$\pm$2.1& $\bm{9.5\pm1.6}$\\
% Inner Ear L & 40.4$\pm$11.8& $\bm{42.3\pm9.9}$ & 59.7$\pm$11.7& $\bm{62.1\pm9.7}$ & 9.4$\pm$1.6 & $\bm{7.6\pm1.5}$ & 7.6$\pm$2.1 & $\bm{5.9\pm1.8}$\\
% Inner Ear R & $\bm{48.9\pm9.5}$ & 45.9$\pm$11.2 & $\bm{67.3\pm10.4}$ & 66.9$\pm$6.1& 8.3$\pm$1.6& $\bm{7.0\pm1.6}$ & $\bm{6.8\pm1.9}$ & $\bm{5.4\pm1.2}$\\
% Mid Ear L & $\bm{62.3\pm14.7}$ & 55.8$\pm$10.0 & $\bm{71.0\pm11.7}$ & 67.5$\pm$6.3& 9.8$\pm$3.3 & $\bm{8.9\pm1.7}$ & 6.4$\pm$2.8 & $\bm{6.1\pm1.1}$\\
% Mid Ear R & $\bm{62.4\pm13.1}$ & 53.6$\pm$12.1 & $\bm{73.3\pm7.7}$ & 65.0$\pm$8.9& 9.7$\pm$2.6 & $\bm{9.0\pm1.3}$ & $\bm{6.2\pm1.9}$ & 6.6$\pm$1.2\\
% TM Joint L & 38.3$\pm$10.1& $\bm{39.0\pm10.9}$ & 59.7$\pm$12.8& $\bm{61.2\pm4.9}$ & 10.2$\pm$2.1& $\bm{7.7\pm1.5}$ & 6.8$\pm$1.9 & $\bm{6.0\pm0.7}$\\
% TM Joint R & $\bm{41.5\pm10.0}$ & 38.3$\pm$9.5 & 59.2$\pm$17.8& $\bm{60.0\pm7.7}$ & 9.5$\pm$1.6 & $\bm{8.0\pm1.3}$ & 8.2$\pm$1.9& $\bm{6.3\pm1.1}$\\
% Spinal Cord & 27.7$\pm$8.2& $\bm{34.7\pm6.9}$ & 38.0$\pm$7.7& $\bm{40.3\pm6.1}$ & 15.5$\pm$2.6& $\bm{13.3\pm2.9}$& 10.7$\pm$1.3& $\bm{8.8\pm0.9}$\\
% Mandible L & $\bm{78.0\pm5.2}$ & 66.7$\pm$5.2& $\bm{86.0\pm2.1}$ & 76.8$\pm$4.7& $\bm{11.0\pm2.7}$ & 12.2$\pm$2.5& $\bm{5.5\pm1.7}$ & 9.0$\pm$4.5\\
% Mandible R & $\bm{69.0\pm4.4}$ & 64.7$\pm$4.9& $\bm{81.2\pm2.7}$ & 75.0$\pm$4.5& 27.0$\pm$6.0& $\bm{13.2\pm2.3}$ & 21.6$\pm$5.4& $\bm{9.6\pm3.5}$ \\
% \hline
% Average & 37.4$\pm$9.7 & $\bm{46.0\pm7.4}$ & 50.3$\pm$9.2 & $\bm{57.2\pm5.6}$ & 12.2$\pm$2.6 & $\bm{9.4\pm1.7}$ & 9.9$\pm$2.2 & $\bm{7.6\pm1.7}$ \\ 
% \bottomrule[1.5pt]
% \end{tabular}}
% \end{subtable}}

% \vspace{0.2cm}

% \textcolor{black}{
% \begin{subtable}{1\textwidth}
% \centering
% \caption{\textcolor{black}{WORD}}
% \scalebox{0.95}{
% \begin{tabular}{ccccccccc} \bottomrule[1.5pt]
% \multirow{3}*{Organs} & \multicolumn{4}{c}{DSC $\uparrow$ (mean$\pm$std \%)} & \multicolumn{4}{c}{HD$_{9}$ $\downarrow$ (mean$\pm$std mm)} \\ \cmidrule(lr){2-5} \cmidrule(lr){6-9}
% ~ & \multicolumn{2}{c}{MedLAM} & \multicolumn{2}{c}{Manual Prompt} & \multicolumn{2}{c}{MedLAM} & \multicolumn{2}{c}{Manual Prompt} \\ \cmidrule(lr){2-3} \cmidrule(lr){4-5} \cmidrule(lr){6-7} \cmidrule(lr){8-9}
% ~ & SAM & MedSAM & SAM & MedSAM & SAM & MedSAM & SAM & MedSAM\\\hline
% Liver & $\bm{66.0\pm10.1}$ & 23.8$\pm$8.6 &  $\bm{84.2\pm6.3}$ & 46.6$\pm$14.7 & 53.6$\pm$10.3 & $\bm{51.5\pm8.6}$ &  $\bm{21.8\pm11.9}$ &  36.8$\pm$11.6 \\
% Spleen & $\bm{60.3\pm15.9}$ &  31.5$\pm$13.7 &  $\bm{85.3\pm5.1}$ &  65.0$\pm$6.3 & $\bm{33.9\pm20.3}$ &  34.0$\pm$21.2 & $\bm{10.8\pm6.5}$ & 16.2$\pm$4.9 \\
% Kidney L & $\bm{82.0\pm16.0}$ &  70.9$\pm$18.8 &  $\bm{92.1\pm1.6}$ &  84.1$\pm$5.2 & 19.3$\pm$18.5 &  $\bm{18.5\pm16.8}$ &  $\bm{6.1\pm1.9}$ &  8.2$\pm$1.6 \\
% Kidney R & $\bm{87.9\pm4.1}$ & 76.9$\pm$5.4 &  $\bm{92.9\pm1.6}$ &  86.4$\pm$3.1 &  $\bm{13.4\pm5.0}$ & 13.5$\pm$2.2 &  $\bm{6.0\pm1.6}$ &  7.7$\pm$1.2 \\
% Stomach & 47.8$\pm$14.6 &  $\bm{39.3\pm13.6}$ & 77.1$\pm$10.6 &  $\bm{80.3\pm4.7}$ & $\bm{50.6\pm25.9}$ &  52.8$\pm$27.2 & 12.7$\pm$4.2 & $\bm{11.5\pm3.6}$ \\
% Gallbladder & 14.0$\pm$15.1 &  $\bm{10.7\pm11.8}$ & $\bm{72.7\pm11.1}$ &  68.8$\pm$7.8 & $\bm{39.5\pm21.4}$ &  38.9$\pm$20.7 &  $\bm{6.0\pm1.8}$ &  6.4$\pm$1.6 \\
% Esophagus & 28.0$\pm$13.6 &  $\bm{23.1\pm10.8}$ &  $\bm{67.0\pm6.4}$ &  63.1$\pm$7.7 & $\bm{24.3\pm12.4}$ &  21.7$\pm$12.4 &  6.9$\pm$4.0 & $\bm{6.7\pm3.2}$ \\
% Pancreas & 32.6$\pm$15.1 & $\bm{20.7\pm10.8}$ &  $\bm{64.4\pm7.7}$ & 46.9$\pm$11.6 & $\bm{35.1\pm12.2}$ &  34.3$\pm$11.8 & 15.9$\pm$6.4 & $\bm{18.3\pm5.5}$ \\
% Duodenum & 23.7$\pm$10.2 & $\bm{18.5\pm5.6}$ & $\bm{54.1\pm13.2}$ & 51.0$\pm$11.9 &  $\bm{34.4\pm9.6}$ &  32.7$\pm$10.0 & $\bm{16.0\pm4.6}$ & 17.1$\pm$5.9 \\
% Colon & 25.6$\pm$9.6 & $\bm{26.6\pm9.0}$ & 41.8$\pm$6.8 & $\bm{44.1\pm8.8}$ & $\bm{57.7\pm17.8}$ & 60.6$\pm$17.7 & $\bm{49.4\pm15.3}$ & 49.1$\pm$13.9 \\
% Intestine & $\bm{37.5\pm7.6}$ & 34.0$\pm$8.6 & $\bm{61.4\pm6.9}$ & 52.5$\pm$8.0 & $\bm{47.2\pm10.2}$ & 44.0$\pm$10.1 & 21.7$\pm$4.2 & $\bm{21.0\pm4.0}$ \\
% Adrenal & 1.7$\pm$2.0 & $\bm{8.3\pm4.5}$ & 17.4$\pm$8.8 & $\bm{26.5\pm5.9}$ & $\bm{35.0\pm9.4}$ & 33.0$\pm$8.6 & $\bm{23.4\pm4.3}$ & 22.6$\pm$2.7 \\
% Rectum & $\bm{48.9\pm15.4}$ & 46.0$\pm$15.6 & 75.5$\pm$4.0 & $\bm{80.0\pm3.6}$ & $\bm{21.2\pm7.4}$ & 20.0$\pm$8.5 & $\bm{7.6\pm0.9}$ & 4.8$\pm$0.6 \\
% Bladder & $\bm{60.2\pm25.3}$ & 57.7$\pm$21.6 & $\bm{83.0\pm15.5}$ & 82.9$\pm$8.0 & $\bm{23.6\pm9.3}$ & 21.6$\pm$8.4 & 7.5$\pm$3.5 & $\bm{6.9\pm2.0}$ \\
% Head of Femur L & $\bm{78.8\pm3.3}$ & 63.0$\pm$4.2 & $\bm{90.5\pm2.7}$ & 80.3$\pm$2.8 & 19.5$\pm$2.4 & $\bm{18.9\pm1.9}$ & $\bm{9.2\pm4.7}$ & 12.3$\pm$2.0 \\
% Head of Femur R & $\bm{76.8\pm4.2}$ & 66.2$\pm$4.1 & $\bm{89.1\pm3.4}$ & 83.2$\pm$2.7 & 21.6$\pm$1.9 & $\bm{18.6\pm1.6}$ & 12.6$\pm$6.1 & $\bm{12.3\pm7.7}$ \\
% \hline
% Average & $\bm{48.2\pm11.4}$ & 38.6$\pm$10.4 & $\bm{71.8\pm7.0}$ & 65.1$\pm$7.0 & 33.1$\pm$12.1 & $\bm{32.2\pm11.7}$ & $\bm{14.6\pm5.1}$ & 16.1$\pm$4.5 \\
% \bottomrule[1.5pt]
% \end{tabular}}
% \end{subtable}
% }
% \label{tab:4}
% \end{table*}

After validating the localization performance of MedLAM, we proceeded to examine the segmentation performance of the proposed MedLSAM framework. MedLSAM utilizes the localization information from MedLAM, together with either SAM or MedSAM, to perform segmentation tasks.

In our evaluation, we also included manually-annotated bounding boxes to serve as the performance \textbf{upper bound} for SAM and MedSAM on the WORD and StructSeg HaN datasets. These were simulated based on the ground-truth masks, following the same approach used for generating bboxes in MedSAM. During inference, a bbox prompt was generated from each ground-truth mask, with a random perturbation of 0-20 pixels introduced to mimic the potential inaccuracies that would be present in manually-drawn bboxes.

It can be observed that MedLSAM demonstrates performance closely approximating that achieved with manual prompts for many organs. For instance, in the case of the brain stem, MedLSAM attains a DSC score of 73.3\%, while the manual prompt reaches 75.9\%. However, for certain abdominal organs with substantial shape variability, such as the pancreas, the performance gap between MedLSAM and manual prompts is more significant, with a DSC score difference reaching up to 35\%. Fig. \ref{fig:seg_result} visually presents the results of using MedLSAM for localization and segmentation tasks on both datasets.
\par Overall, MedLSAM performs more comparably to manual prompts in the case of head and neck organs, representing a pioneering and effective attempt towards fully automating the use of SAM in medical imaging.

\subsubsection{Ablation study of SPL}
To validate the effectiveness of Sub-Patch Localization (SPL), we conducted experiments on the StructSeg HaN dataset, comparing the performance of SPL at different sub-patch intervals $n$ mm against the Whole-Patch Localization (WPL) approach. As described in Sec. \ref{sec:4.2}, the WPL method relies on the extreme coordinates of the entire organ for localization. On the other hand, SPL employs a more granular approach by segmenting the organ into multiple sub-patches at intervals of $n$ mm.
\par As shown in Table \ref{tab:5}, it is evident that the SPL approach consistently outperforms WPL. Notably, for organs with significant morphological variations, such as the brain stem and spinal cord, the DSC score improves by more than 7\% and 20\%, respectively.  Moreover, when $n$ is relatively small, SPL enables finer-grained localization by concentrating on smaller, more specific portions of the organ.  This results in bounding boxes that better approximate the organ's actual shape at each slice, thereby enhancing the accuracy of the subsequent segmentation process. 

\section{Discussion \& Conclusions}
In this study, we have introduced MedLSAM, the first completely automated, two-stage medical adaptation of the SAM model, designed to significantly alleviate the annotation workload in the segmentation of medical images. The first stage, MedLAM, serves as a pioneering localization foundation model, capable of precisely identifying anatomical structures in 3D medical images. The second stage leverages the strength of SAM and its medical adaptations for meticulous image segmentation. Together, these components form a fully automated pipeline that significantly reduces the reliance on manual annotations and annotated data.

This endeavor was fueled by the observation that the spatial distributions of organs across different patients maintain strong similarities. Consequently, MedLAM was designed to project every part of the scans onto a shared 3D latent coordinate system, accurately localizing target anatomical parts within the body. Coupling this approach with SAM's segmentation capabilities led to an efficient and accurate process for image segmentation.

Moreover, MedLSAM demonstrated its effectiveness across two 3D datasets covering 38 different organs, providing robust evidence of its versatility. Importantly, this automated approach is not burdened by an increased annotation workload as data size increases. It also holds promise for direct integration with potential future 3D SAM models in the medical field, which could further enhance its performance and utility.

\section*{Acknowledgments}

\bibliographystyle{model2-names.bst}\biboptions{authoryear}
\bibliography{medlsam.bib}

\end{document}

%%
